\makeatletter
\expandafter\gdef\csname CGTip@sec:introduction\endcsname{Section 1: Introduction (p. 3)}
\expandafter\gdef\csname CGTip@subsec:clock_intro\endcsname{Section 1.1: The Riemann--von Mangoldt clock and the ordinates (p. 3)}
\expandafter\gdef\csname CGTip@eq:intro_rvm\endcsname{Eq. (1.1), p. 4: N(T)=N\_0(T)+(1)/(pi)arg zeta ((1)/(2)+iT)+O((1)/(T)), N\_0(T)=(T)/(2pi)log(T)/(2pi e)+(7)/(8),}
\expandafter\gdef\csname CGTip@eq:intro_clock\endcsname{Eq. (1.2), p. 4: N\_0(tau\_k)=k-(1)/(2), tau\_k=(2pi (k-(11)/(8)))/(W\_0((k-(11)/(8))/e))(k>=1),}
\expandafter\gdef\csname CGTip@eq:intro_defect\endcsname{Eq. (1.3), p. 4: s\_k=k-(1)/(2)-N\_0(gamma\_k), u\_k(X)=s\_k-S\_X(gamma\_k), S\_X(t)=-(1)/(pi)sum\_p<= Xarg (1-p\textasciicircum{}-1/2-it).}
\expandafter\gdef\csname CGTip@subsec:known_intro\endcsname{Section 1.2: What is known and only used (p. 4)}
\expandafter\gdef\csname CGTip@subsec:main_results\endcsname{Section 1.3: What this paper adds (p. 5)}
\expandafter\gdef\csname CGTip@eq:intro_main\endcsname{Eq. (1.4), p. 5: (1)/(no.I\_T)sum\_kin I\_Tu\_k(X)\textasciicircum{}2=(1)/(2pi\textasciicircum{}2)log(1)/(c)+O(1+sqrt(log(1/c))),}
\expandafter\gdef\csname CGTip@eq:intro_D\endcsname{Eq. (1.5), p. 5: (1)/(no.I\_T)sum\_kin I\_T(tau\_k\textasciicircum{}X-gamma\_k)\textasciicircum{}2=(log log log T+O(sqrt(log log log T)))/(log\textasciicircum{}2T).}
\expandafter\gdef\csname CGTip@eq:intro_E\endcsname{Eq. (1.6), p. 6: limsup\_k->inf|tau\_k-gamma\_k|<=(1)/(log 2)int\_0\textasciicircum{}1(e(sigma) dsigma)/(sqrt(sigma(2-sigma)))<0.60832}
\expandafter\gdef\csname CGTip@eq:intro_ratio\endcsname{Eq. (1.7), p. 6: (<(tau\_k\textasciicircum{}X-gamma\_k)\textasciicircum{}2>\_I\_T)/(<(tau\_k-gamma\_k)\textasciicircum{}2>\_I\_T)\textasciitilde{}(log log log T)/(2log log T)->0,}
\expandafter\gdef\csname CGTip@sec:notation\endcsname{Section 2: Notation (p. 7)}
\expandafter\gdef\csname CGTip@subsec:symbols\endcsname{Section 2.1: List of symbols (p. 7)}
\expandafter\gdef\csname CGTip@subsec:counting\endcsname{Section 2.2: Counting functions (p. 9)}
\expandafter\gdef\csname CGTip@eq:RvM\endcsname{Eq. (2.1), p. 9: N(t)=(1)/(pi)vartheta(t)+1+S(t), S(t)=(1)/(pi)arg zeta ((1)/(2)+it)=O(log t),}
\expandafter\gdef\csname CGTip@eq:N0_def\endcsname{Eq. (2.2), p. 9: Theta\_0(t)=(t)/(2)log(t)/(2pi e)-(pi)/(8), N\_0(t)=(1)/(pi)Theta\_0(t)+1, varpi(t)=vartheta(t)-Theta\_0(t).}
\expandafter\gdef\csname CGTip@eq:N_decomposition\endcsname{Eq. (2.3), p. 9: N(t)=N\_0(t)+S(t)+(varpi(t))/(pi), L(t)=N\_0'(t)=(1)/(2pi)log(t)/(2pi).}
\expandafter\gdef\csname CGTip@eq:IT_size\endcsname{Eq. (2.4), p. 9: no.I\_T=TL(T)+O(T).}
\expandafter\gdef\csname CGTip@subsec:clock\endcsname{Section 2.3: The RvM clock and the displacements (p. 10)}
\expandafter\gdef\csname CGTip@eq:clock\endcsname{Eq. (2.5), p. 10: tau\_x=(2pi (x-(11)/(8)))/(W\_0((x-(11)/(8))/e)), (dtau\_x)/(dx)=(1)/(L(tau\_x))=(2pi)/(log(tau\_x/2pi)),}
\expandafter\gdef\csname CGTip@def:displacements\endcsname{Definition 2.1 (p. 10): For k>=1, s\_k=k-(1)/(2)-N\_0(gamma\_k) and delta\_k=tau\_k-gamma\_k.}
\expandafter\gdef\csname CGTip@eq:s_delta\endcsname{Eq. (2.6), p. 10: s\_k=L(xi\_k) delta\_kfor some xi\_k between gamma\_k and tau\_k .}
\expandafter\gdef\csname CGTip@lem:s_bounds\endcsname{Lemma 2.2 (p. 10): For every k>=1, N(gamma\_k\textasciicircum{}-)-N\_0(gamma\_k)+(1)/(2)<= s\_k<= N(gamma\_k\textasciicircum{}+)-N\_0(gamma\_k)-(1)/(2), and |s\_k-(varpi(gamma\_k))/(pi)|<=max (|S(gamma\_k\textasciicircum{}-)|,|S(gamma\_k\textasciicircum{}+)|)-(1)/(2) .}
\expandafter\gdef\csname CGTip@eq:s_bounds\endcsname{Eq. (2.7), p. 10: N(gamma\_k\textasciicircum{}-)-N\_0(gamma\_k)+(1)/(2)<= s\_k<= N(gamma\_k\textasciicircum{}+)-N\_0(gamma\_k)-(1)/(2),}
\expandafter\gdef\csname CGTip@eq:s_S_bound\endcsname{Eq. (2.8), p. 10: |s\_k-(varpi(gamma\_k))/(pi)|<=max (|S(gamma\_k\textasciicircum{}-)|,|S(gamma\_k\textasciicircum{}+)|)-(1)/(2) .}
\expandafter\gdef\csname CGTip@subsec:euler_phase\endcsname{Section 2.4: Euler phases and the corrected clock (p. 10)}
\expandafter\gdef\csname CGTip@eq:SA_def\endcsname{Eq. (2.9), p. 10: S\_A(t)=sum\_nin Aa\_nsin(tlog n), v\_A=(1)/(2)sum\_nin Aa\_n\textasciicircum{}2 ;}
\expandafter\gdef\csname CGTip@eq:SX_def\endcsname{Eq. (2.10), p. 10: S\_X(t)=-(1)/(pi)sum\_p<= Xarg (1-p\textasciicircum{}-1/2-it)=-(1)/(pi)sum\_p<= Xsum\_m>=1(sin(mtlog p))/(m p\textasciicircum{}m/2)=S\_A\_X(t).}
\expandafter\gdef\csname CGTip@eq:sigma_M_def\endcsname{Eq. (2.11), p. 11: |S\_X(t)|<=sigma\_X:=(1)/(pi)sum\_p<= Xsum\_m>=1(1)/(m p\textasciicircum{}m/2), |S\_X'(t)|<= M\_X:=(1)/(pi)sum\_p<= X(log p)/(sqrt(p)-1),}
\expandafter\gdef\csname CGTip@def:defect\endcsname{Definition 2.3 (p. 11): For X>=2 let F\_X=N\_0+S\_X and u\_k(X)=k-(1)/(2)-F\_X(gamma\_k)=s\_k-S\_X(gamma\_k). The one-step clock and its residual are tau\_k\textasciicircum{}X=tau\_k-(S\_X(tau\_k))/(L(tau\_k)), delta\_k\textasciicircum{}X=tau\_k\textasciicircum{}X-gamma\_k, r\_k(X)=L(tau\_k) delta\_k\textasciicircum{}X .}
\expandafter\gdef\csname CGTip@eq:one_step_clock\endcsname{Eq. (2.12), p. 11: tau\_k\textasciicircum{}X=tau\_k-(S\_X(tau\_k))/(L(tau\_k)), delta\_k\textasciicircum{}X=tau\_k\textasciicircum{}X-gamma\_k, r\_k(X)=L(tau\_k) delta\_k\textasciicircum{}X .}
\expandafter\gdef\csname CGTip@sec:known_results\endcsname{Section 3: External results (p. 11)}
\expandafter\gdef\csname CGTip@thm:known_basic\endcsname{Theorem 3.1 ([ref, Thms. 9.4, 9.9 and S 9.12], [ref]) (p. 11): S(t)=O(log t); gamma\_k+1-gamma\_k is bounded; and S\_1(t)=int\_0\textasciicircum{}tS(y) dy=O(log t).}
\expandafter\gdef\csname CGTip@thm:known_BW\endcsname{Theorem 3.2 (Bellotti--Wong [ref, Thm. 1.1, Cor. 1.5 and (1.2)]) (p. 11): For T>= e, |N(T)-(T)/(2pi)log(T)/(2pi e)|<=BW(T) and |S(T)|<=BWS(T), where BWS(t)=0.10076log t+min 0.24460log log t+7.20844, 1.68845log log t+1.50956. Moreover |S(T)|<=2.5167 for 0<= T<=30 610 046 000.}
\expandafter\gdef\csname CGTip@eq:BWS_def\endcsname{Eq. (3.1), p. 11: BWS(t)=0.10076log t+min 0.24460log log t+7.20844, 1.68845log log t+1.50956.}
\expandafter\gdef\csname CGTip@thm:known_selberg_moments\endcsname{Theorem 3.3 (Selberg [ref]; see [ref, p. 225]) (p. 11): For every integer j>=1, int\_T\textasciicircum{}2TS(t)\textasciicircum{}2j dt=kappa\_j T(log log T)\textasciicircum{}j+O\_j(T(log log T)\textasciicircum{}j-1/2).}
\expandafter\gdef\csname CGTip@thm:known_korolev\endcsname{Theorem 3.4 (Selberg; Korolev [ref], [ref, Thm. 4.2]) (p. 11): Let Delta\_k be the integer such that g\_k+Delta\_k-2<gamma\_k<= g\_k+Delta\_k-1 (in [ref] the Gram points are t\_m=g\_m-1 and t\_m-1<gamma\_k<= t\_m with m=k+Delta\_k), and let 0<varepsilon\_0<10\textasciicircum{}-3. For every fixed j>=1, as n\_0->inf, uniformly for n\_0\textasciicircum{}27/82+0.9varepsilon\_0<= h\_0<= n\_0\textasciicircum{}27/82+varepsilon\_0, sum\_n\_0<k<= n\_0+h\_0Del...}
\expandafter\gdef\csname CGTip@eq:korolev_moments\endcsname{Eq. (3.2), p. 11: sum\_n\_0<k<= n\_0+h\_0Delta\_k\textasciicircum{}2j=kappa\_jh\_0(log log n\_0)\textasciicircum{}j(1+o(1)), |sum\_n\_0<k<= n\_0+h\_0Delta\_k\textasciicircum{}2j-1|<<\_j h\_0(log log n\_0)\textasciicircum{}j-1.}
\expandafter\gdef\csname CGTip@thm:known_selberg_approx\endcsname{Theorem 3.5 (Selberg's approximation in the form of Fujii [ref, S 3, pp. 232--233]) (p. 12): For every integer k>=1, int\_T\textasciicircum{}2T(S(t)-S\textasciicircum{}rm pr\_Y\_k(t))\textasciicircum{}2kdt<=(Ck\textasciicircum{}2)\textasciicircum{}2k T with Y\_k=T\textasciicircum{}1/(80k) and an absolute constant C. In particular int\_T\textasciicircum{}2T(S-S\textasciicircum{}rm pr\_Y\_0)\textasciicircum{}4<< T with Y\_0=Y\_2=T\textasciicircum{}1/160.}
\expandafter\gdef\csname CGTip@thm:known_fujii_difference\endcsname{Theorem 3.6 (Fujii [ref, Main Theorem, p. 225]) (p. 12): For every integer k>=1 there is C\_k such that, for T>= T\_0 and 0<h<=1, [ int\_T\textasciicircum{}2T(S(t+h)-S(t))\textasciicircum{}2kdt<= C\_k T(log(2+hlog T))\textasciicircum{}k . ]}
\expandafter\gdef\csname CGTip@thm:known_MV\endcsname{Theorem 3.7 (Montgomery--Vaughan [ref, Cor. 3]) (p. 12): For complex z\_n, int\_T\textasciicircum{}2T|sum\_n<= Yz\_nn\textasciicircum{}-it|\textasciicircum{}2dt=sum\_n<= Y|z\_n|\textasciicircum{}2(T+O(n)) with an absolute implied constant.}
\expandafter\gdef\csname CGTip@thm:known_mertens\endcsname{Theorem 3.8 (Mertens [ref]; see [ref, Ch. I.1, Thms. 1.8 and 1.10]) (p. 12): For y>=2, sum\_p<= y(log p)/(p)=log y+O(1) and sum\_p<= y(1)/(p)=log log y+O(1).}
\expandafter\gdef\csname CGTip@thm:known_KL\endcsname{Theorem 3.9 (Kusmin--Landau [ref, Thm. 2.1]) (p. 12): Let Omega be real with continuous monotone derivative on [y\_1,y\_2] and ||Omega'||>=1/H\_rm KL>0 there, where ||*|| is the distance to the nearest integer. Then sum\_y\_1<n<= y\_2e\textasciicircum{}2pi iOmega(n)<< H\_rm KL.}
\expandafter\gdef\csname CGTip@thm:known_stationary\endcsname{Theorem 3.10 (Stationary phase [ref, Lemma 4.6]) (p. 12): Let F be real with three derivatives on [a,b], 0<lambda\_2<=-F''<= Alambda\_2 and |F'''|<= Alambda\_3 there, and F'(c)=0 with a<= c<= b. Then int\_a\textasciicircum{}be\textasciicircum{}iF(x)dx= ((2pi)\textasciicircum{}1/2e\textasciicircum{}iF(c)-ipi/4)/(|F''(c)|\textasciicircum{}1/2)+O(lambda\_2\textasciicircum{}-4/5lambda\_3\textasciicircum{}1/5); +O(min ((1)/(|F'(a)|),lambda\_2\textasciicircum{}-1/2))+O(min ((1)/(|F'(b)|),lambda\_2\textasciicircum{}-1/2)), and if F' does...}
\expandafter\gdef\csname CGTip@thm:known_vdC\endcsname{Theorem 3.11 (Van der Corput's exponential sum to integral [ref, Lemma 4.7]) (p. 12): Let f be real with a continuous, steadily decreasing derivative on (a,b), f'(b)=alpha, f'(a)=beta. Then for any 0<eta<1, [ sum\_a<n<= be\textasciicircum{}2pi if(n)=sum\_alpha-eta<nu<beta+etaint\_a\textasciicircum{}be\textasciicircum{}2pi i(f(x)-nu x)dx+O(log(beta-alpha+2)). ]}
\expandafter\gdef\csname CGTip@thm:known_vdC_weighted\endcsname{Theorem 3.12 (Van der Corput's exponential sum to integral, weighted form [ref, Lemma 4.10]) (p. 13): Let f be real with a continuous, steadily decreasing derivative on (a,b), f'(b)=alpha, f'(a)=beta, and let g>0 be decreasing with continuous derivative such that | g'| is steadily decreasing. Then for any 0<eta<1, [ sum\_a<n<= b g(n)e\textasciicircum{}2pi if(n)=sum\_alpha-eta<nu<beta+etaint\_a\textasciicircum{}b g(x)e\textasciicircum{}2pi i(f(x)-nu x)dx+O(g(a)log(beta-...}
\expandafter\gdef\csname CGTip@thm:known_vdC2\endcsname{Theorem 3.13 (Van der Corput's second-derivative test [ref, Thm. 5.9]) (p. 13): Let f be real and twice differentiable on [a,b], b>= a+1, with 0<lambda\_2<= f''<= hlambda\_2 (or lambda\_2<=-f''<= hlambda\_2) throughout. Then sum\_a<n<= be\textasciicircum{}2pi if(n)=O(h(b-a)lambda\_2\textasciicircum{}1/2)+O(lambda\_2\textasciicircum{}-1/2).}
\expandafter\gdef\csname CGTip@thm:known_landau\endcsname{Theorem 3.14 (Landau [ref]; uniform version Gonek [ref], as stated in [ref, (1.1)]) (p. 13): For integers n>=2 and real t>=2, sum\_0<gamma <= tn\textasciicircum{}rho=-(t)/(2pi)Lambda(n)+O(nlog(2nt)log log(3n)), zeros counted with multiplicity, with an absolute implied constant.}
\expandafter\gdef\csname CGTip@thm:known_density\endcsname{Theorem 3.15 (Selberg [ref, Thm. 9.19(C)]) (p. 13): Let N\textasciicircum{}+(sigma,t) be the number of zeros beta+igamma of the Riemann zeta function with beta>sigma and 0<gamma <= t, counted with multiplicity. Uniformly for (1)/(2)<=sigma <=1 and t>=2, N\textasciicircum{}+(sigma,t)<< t\textasciicircum{}1-(1)/(4)(sigma-(1)/(2))log t.}
\expandafter\gdef\csname CGTip@thm:known_selberg_clt\endcsname{Theorem 3.16 (Selberg [ref]; Tsang [ref]; see [ref, (1.1)]) (p. 13): For fixed real y\_1<y\_2, (1)/(T)meastin[T,2T]:y\_1<S(t)/ s\_T<= y\_2->(1)/(sqrt(2pi))int\_y\_1\textasciicircum{}y\_2e\textasciicircum{}-y\textasciicircum{}2/2dy, where s\_T=((1)/(2pi\textasciicircum{}2)log log T)\textasciicircum{}1/2.}
\expandafter\gdef\csname CGTip@thm:known_maximal\endcsname{Theorem 3.17 (Hardy--Littlewood [ref, Thm. I.1], as quoted in [ref, (1.2)]; Kinnunen [ref, Thm. 1.4]) (p. 13): Let Mf(x)=sup\_y>0(1)/(2y)int\_x-y\textasciicircum{}x+y|f| be the centred maximal function of fin L\textasciicircum{}1\_rm loc(R). For 1<p<inf and fin L\textasciicircum{}p(R), || Mf||\_p<= A\_p||f||\_p, where A\_p is bounded for p>=2. If f and f' belong to L\textasciicircum{}p(R), 1<p<inf, then Mf has a locally absolutely continuous representative, (Mf)'in L\textasciicircum{}p(R), and |(Mf)'|<= M(f') almos...}
\expandafter\gdef\csname CGTip@thm:known_carleson\endcsname{Theorem 3.18 (Carleson [ref], Hunt [ref]; on R see [ref]) (p. 13): For fin L\textasciicircum{}2(R)n L\textasciicircum{}4(R) with Fourier transform Ff(omega)=int\_R f(t)e\textasciicircum{}-iomega tdt, let Cf(x)=sup\_omega\_1in R|(1)/(2pi)int\_-inf\textasciicircum{}omega\_1 Ff(omega)e\textasciicircum{}iomega xdomega |. Then || Cf||\_4<<||f||\_4.}
\expandafter\gdef\csname CGTip@thm:known_exponents\endcsname{Theorem 3.19 ([ref, Lemma 2.8 and (7.57)]; Bourgain [ref, S 5, Thm. 6]) (p. 13): (a) If (k, l) is an exponent pair, so is A(k, l)=((k)/(2 k+2),(1)/(2)+(l)/(2 k+2)). (b) If (k, l) is an exponent pair with l- k>=(1)/(2), then zeta(l- k+it)<<|t|\textasciicircum{} klog|t| for |t|>=2. (c) For every varepsilon>0, ((13)/(84)+varepsilon,(55)/(84)+varepsilon) is an exponent pair.}
\expandafter\gdef\csname CGTip@thm:known_PL\endcsname{Theorem 3.20 (Phragmen--Lindelof principle in a strip [ref, Lemma 2.5], [ref, Lemma 3]) (p. 14): Let sigma\_1<sigma\_2 and e\_1>= e\_2>=0. If |zeta(sigma\_i+it)|<= C\_i|t|\textasciicircum{}e\_ilog|t| for |t|>=2 and i=1,2, then |zeta(sigma+it)|<<|t|\textasciicircum{}e(sigma)log|t| uniformly for sigma\_1<=sigma <=sigma\_2 and |t|>=2, where e is the linear function with e(sigma\_i)=e\_i.}
\expandafter\gdef\csname CGTip@thm:known_backlund\endcsname{Theorem 3.21 (Backlund's method [ref, Lemmas 2.1, 2.3 and 2.4], after [ref]) (p. 14): Let T be large, varkappa in N, f\_varkappa(s)=(1)/(2)[((s+iT-1)zeta(s+iT))\textasciicircum{}varkappa+((s-iT-1)zeta(s-iT))\textasciicircum{}varkappa ], and for a disc D(sigma\_o,varrho) let Z\_varkappa(sigma\_o,varrho)=(1)/(varkappa)sum\_zlog(varrho)/(|z-sigma\_o|) over the zeros z of f\_varkappa in D(sigma\_o,varrho). Let 1<sigma\_o<varrho, sigma\textasciicircum{}sharp=sigma...}
\expandafter\gdef\csname CGTip@rem:korolev\endcsname{Remark 3.22 (Known consequences for the clock) (p. 14): (a) Moments and law of s\_k. Since vartheta(gamma\_k)/pi in(k+Delta\_k-2,k+Delta\_k-1] and vartheta/pi=N\_0-1+varpi/pi, one has s\_k+Delta\_k=k+Delta\_k-(1)/(2)-N\_0(gamma\_k)in[-(1)/(2)+(varpi(gamma\_k))/(pi),(1)/(2)+(varpi(gamma\_k))/(pi)), so s\_k=-Delta\_k+ d\_k with | d\_k|<1. The index set I\_T is the integer interval (N(T),N(...}
\expandafter\gdef\csname CGTip@eq:s_moments\endcsname{Eq. (3.3), p. 14: < s\_k\textasciicircum{}2j>\_I\_T=kappa\_j(log log T)\textasciicircum{}j(1+o(1)), < s\_k\textasciicircum{}2j-1>\_I\_T=O\_j((log log T)\textasciicircum{}j-1).}
\expandafter\gdef\csname CGTip@eq:BW_pointwise\endcsname{Eq. (3.4), p. 15: -(BW(gamma\_k)+3/8)/(L(tau\_k))<=delta\_k<=(BW(gamma\_k)-11/8)/(L(gamma\_k))(k>=1).}
\expandafter\gdef\csname CGTip@sec:identities\endcsname{Section 4: Theorem A: an exact defect identity at the zeros (p. 15)}
\expandafter\gdef\csname CGTip@thm:identity_general\endcsname{Theorem 4.1 (p. 15): Let 1<= k\_1<k\_2 be indices, F a real C\textasciicircum{}1 function on an open interval containing [gamma\_k\_1,gamma\_k\_2], and u\_k=k-(1)/(2)-F(gamma\_k) for k\_1<= k<= k\_2. For every integer j>=1, int\_gamma\_k\_1\textasciicircum{}gamma\_k\_2(N(t)-F(t))\textasciicircum{}2jF'(t) dt=sum\_k=k\_1+1\textasciicircum{}k\_2-1Psi\_j(u\_k)+(((1)/(2)+u\_k\_1)\textasciicircum{}2j+1+((1)/(2)-u\_k\_2)\textasciicircum{}2j+1)/(2j+1), where Psi\_j(u)=...}
\expandafter\gdef\csname CGTip@eq:identity_general\endcsname{Eq. (4.1), p. 15: int\_gamma\_k\_1\textasciicircum{}gamma\_k\_2(N(t)-F(t))\textasciicircum{}2jF'(t) dt=sum\_k=k\_1+1\textasciicircum{}k\_2-1Psi\_j(u\_k)+(((1)/(2)+u\_k\_1)\textasciicircum{}2j+1+((1)/(2)-u\_k\_2)\textasciicircum{}2j+1)/(2j+1),}
\expandafter\gdef\csname CGTip@eq:Psi_def\endcsname{Eq. (4.2), p. 15: Psi\_j(u)=(((1)/(2)+u)\textasciicircum{}2j+1+((1)/(2)-u)\textasciicircum{}2j+1)/(2j+1)=sum\_l=0\textasciicircum{}jbinom(2j+1,2l)(4\textasciicircum{}l-j)/(2j+1) u\textasciicircum{}2l=u\textasciicircum{}2j+sum\_l<jbinom(2j+1,2l)(4\textasciicircum{}l-j)/(2j+1)u\textasciicircum{}2l.}
\expandafter\gdef\csname CGTip@eq:identity_quadratic\endcsname{Eq. (4.3), p. 15: sum\_k=k\_1+1\textasciicircum{}k\_2-1u\_k\textasciicircum{}2=int\_gamma\_k\_1\textasciicircum{}gamma\_k\_2(N(t)-F(t))\textasciicircum{}2F'(t) dt-(k\_2-k\_1-1)/(12)-(((1)/(2)+u\_k\_1)\textasciicircum{}3+((1)/(2)-u\_k\_2)\textasciicircum{}3)/(3) .}
\expandafter\gdef\csname CGTip@eq:identity_one_gap\endcsname{Eq. (4.4), p. 15: int\_gamma\_k\textasciicircum{}gamma\_k+1(k-F(t))\textasciicircum{}2jF'(t) dt =((k-F(gamma\_k))\textasciicircum{}2j+1-(k-F(gamma\_k+1))\textasciicircum{}2j+1)/(2j+1); =(((1)/(2)+u\_k)\textasciicircum{}2j+1+((1)/(2)-u\_k+1)\textasciicircum{}2j+1)/(2j+1),}
\expandafter\gdef\csname CGTip@thm:identity_g\endcsname{Theorem 4.2 (p. 16): Let k\_1<k\_2, F and u\_k be as in Theorem 4.1, and let g:R -> C be continuous with a primitive G. Then int\_gamma\_k\_1\textasciicircum{}gamma\_k\_2 g(N(t)-F(t))F'(t) dt=sum\_k=k\_1+1\textasciicircum{}k\_2-1[ G(u\_k+(1)/(2))- G(u\_k-(1)/(2))]+ G(u\_k\_1+(1)/(2))- G(u\_k\_2-(1)/(2)).}
\expandafter\gdef\csname CGTip@eq:identity_g\endcsname{Eq. (4.5), p. 16: int\_gamma\_k\_1\textasciicircum{}gamma\_k\_2 g(N(t)-F(t))F'(t) dt=sum\_k=k\_1+1\textasciicircum{}k\_2-1[ G(u\_k+(1)/(2))- G(u\_k-(1)/(2))]+ G(u\_k\_1+(1)/(2))- G(u\_k\_2-(1)/(2)).}
\expandafter\gdef\csname CGTip@cor:identity_exp\endcsname{Corollary 4.3 (p. 16): For real omega !=0, (2sin(omega/2))/(omega)sum\_k=k\_1+1\textasciicircum{}k\_2-1e\textasciicircum{}iomega u\_k=int\_gamma\_k\_1\textasciicircum{}gamma\_k\_2e\textasciicircum{}iomega(N(t)-F(t))F'(t) dt-(e\textasciicircum{}iomega(u\_k\_1+1/2)-e\textasciicircum{}iomega(u\_k\_2-1/2))/(iomega).}
\expandafter\gdef\csname CGTip@eq:identity_exp\endcsname{Eq. (4.6), p. 16: (2sin(omega/2))/(omega)sum\_k=k\_1+1\textasciicircum{}k\_2-1e\textasciicircum{}iomega u\_k=int\_gamma\_k\_1\textasciicircum{}gamma\_k\_2e\textasciicircum{}iomega(N(t)-F(t))F'(t) dt-(e\textasciicircum{}iomega(u\_k\_1+1/2)-e\textasciicircum{}iomega(u\_k\_2-1/2))/(iomega).}
\expandafter\gdef\csname CGTip@cor:signed_law\endcsname{Corollary 4.4 (Signed law of s\_k from the t-law of S) (p. 16): For fixed real eta, < e\textasciicircum{}ieta s\_k/ s\_T>\_I\_T-> e\textasciicircum{}-eta\textasciicircum{}2/2. Hence the empirical law of s\_k/ s\_T on I\_T tends to the standard normal law.}
\expandafter\gdef\csname CGTip@rem:twelfth\endcsname{Remark 4.5 (p. 16): Psi\_j(u)=int\_u-1/2\textasciicircum{}u+1/2x\textasciicircum{}2jdx is the mean of x\textasciicircum{}2j over the unit interval centred at u; in particular Psi\_1(u)=u\textasciicircum{}2+(1)/(12) and (1)/(12)=int\_-1/2\textasciicircum{}1/2x\textasciicircum{}2dx. The constant comes from the unit jump of N-F at each zero, through the primitive in Theorem 4.2; when F is not monotone, N-F need not run through an interval of...}
\expandafter\gdef\csname CGTip@cor:identity_clock\endcsname{Corollary 4.6 (p. 17): For X>=2 and 1<= k\_1<k\_2, sum\_k=k\_1+1\textasciicircum{}k\_2-1u\_k(X)\textasciicircum{}2= int\_gamma\_k\_1\textasciicircum{}gamma\_k\_2(S-S\_X+(varpi)/(pi))\textasciicircum{}2(L+S\_X')dt; -(k\_2-k\_1-1)/(12)-(((1)/(2)+u\_k\_1(X))\textasciicircum{}3+((1)/(2)-u\_k\_2(X))\textasciicircum{}3)/(3) .}
\expandafter\gdef\csname CGTip@eq:identity_LeClair\endcsname{Eq. (4.7), p. 17: sum\_k=k\_1+1\textasciicircum{}k\_2-1u\_k(X)\textasciicircum{}2= int\_gamma\_k\_1\textasciicircum{}gamma\_k\_2(S-S\_X+(varpi)/(pi))\textasciicircum{}2(L+S\_X')dt; -(k\_2-k\_1-1)/(12)-(((1)/(2)+u\_k\_1(X))\textasciicircum{}3+((1)/(2)-u\_k\_2(X))\textasciicircum{}3)/(3) .}
\expandafter\gdef\csname CGTip@lem:first_moment_identity\endcsname{Lemma 4.7 (First-moment identity) (p. 17): Let 2pi<t\_1<t\_2 be points which are not ordinates and psi in C\textasciicircum{}1[t\_1,t\_2], real or complex. Then sum\_t\_1<gamma\_k<= t\_2s\_kpsi(gamma\_k)=int\_t\_1\textasciicircum{}t\_2G(t)psi(t)L(t) dt+(1)/(2)[psi G\textasciicircum{}2]\_t\_1\textasciicircum{}t\_2-(1)/(2)int\_t\_1\textasciicircum{}t\_2G(t)\textasciicircum{}2psi'(t) dt .}
\expandafter\gdef\csname CGTip@eq:first_moment_identity\endcsname{Eq. (4.8), p. 17: sum\_t\_1<gamma\_k<= t\_2s\_kpsi(gamma\_k)=int\_t\_1\textasciicircum{}t\_2G(t)psi(t)L(t) dt+(1)/(2)[psi G\textasciicircum{}2]\_t\_1\textasciicircum{}t\_2-(1)/(2)int\_t\_1\textasciicircum{}t\_2G(t)\textasciicircum{}2psi'(t) dt .}
\expandafter\gdef\csname CGTip@eq:first_moment_measures\endcsname{Eq. (4.9), p. 17: d(G\textasciicircum{}2)=-2GL dt+sum\_t\_1<gamma<t\_22G\textasciicircum{}*(gamma)nu(gamma) delta\_gamma ,}
\expandafter\gdef\csname CGTip@sec:main_theorem\endcsname{Section 5: Theorem B: how much prime information the clock needs (p. 17)}
\expandafter\gdef\csname CGTip@thm:main\endcsname{Theorem 5.1 (p. 17): There are absolute constants T\_0 and C such that for T>= T\_0 and every X with 2<= X<= T\textasciicircum{}1/2, writing c=log X/log T, |(1)/(no.I\_T)sum\_kin I\_Tu\_k(X)\textasciicircum{}2-(1)/(2pi\textasciicircum{}2)log(1)/(c)|<= C(1+sqrt(log(1/c))).}
\expandafter\gdef\csname CGTip@eq:main_theorem\endcsname{Eq. (5.1), p. 17: |(1)/(no.I\_T)sum\_kin I\_Tu\_k(X)\textasciicircum{}2-(1)/(2pi\textasciicircum{}2)log(1)/(c)|<= C(1+sqrt(log(1/c))).}
\expandafter\gdef\csname CGTip@cor:fixed_cutoff\endcsname{Corollary 5.2 (fixed cutoff) (p. 17): For fixed X>=2, (1)/(no.I\_T)sum\_kin I\_Tu\_k(X)\textasciicircum{}2=(1)/(2pi\textasciicircum{}2)log log T+O\_X(sqrt(log log T)), and no.kin I\_T:|u\_k(X)|<= y\_0=o(no.I\_T) for every fixed y\_0>0.}
\expandafter\gdef\csname CGTip@cor:power_cutoff\endcsname{Corollary 5.3 (power cutoff) (p. 18): For fixed 0<c<=(1)/(2) and X=T\textasciicircum{}c, limsup\_T->inf(1)/(no.I\_T)sum\_kin I\_Tu\_k(T\textasciicircum{}c)\textasciicircum{}2<=(1)/(2pi\textasciicircum{}2)log(1)/(c)+C(1+sqrt(log(1/c))), and no.kin I\_T:|u\_k(T\textasciicircum{}c)|>w\_T=o(no.I\_T) whenever w\_T->inf. The same holds for X=N(T)\textasciicircum{}c with fixed 0<c<(1)/(2), since then X=T\textasciicircum{}c' with c'=c+O(log log T/log T).}
\expandafter\gdef\csname CGTip@eq:power_cutoff\endcsname{Eq. (5.2), p. 18: limsup\_T->inf(1)/(no.I\_T)sum\_kin I\_Tu\_k(T\textasciicircum{}c)\textasciicircum{}2<=(1)/(2pi\textasciicircum{}2)log(1)/(c)+C(1+sqrt(log(1/c))),}
\expandafter\gdef\csname CGTip@cor:dichotomy\endcsname{Corollary 5.4 (dichotomy) (p. 18): Let X=X(T) with 2<= X<= T\textasciicircum{}1/2. The mean square of u\_k(X) over I\_T is bounded as T->inf if and only if liminf\_T->inflog X/log T>0. For fixed X\_0>=2 and cin(0,(1)/(2)], (1)/(no.I\_T)sum\_kin I\_T(u\_k(X\_0)\textasciicircum{}2-u\_k(T\textasciicircum{}c)\textasciicircum{}2)=(1)/(2pi\textasciicircum{}2)log log T+O(sqrt(log log T)) .}
\expandafter\gdef\csname CGTip@eq:dichotomy\endcsname{Eq. (5.3), p. 18: (1)/(no.I\_T)sum\_kin I\_T(u\_k(X\_0)\textasciicircum{}2-u\_k(T\textasciicircum{}c)\textasciicircum{}2)=(1)/(2pi\textasciicircum{}2)log log T+O(sqrt(log log T)) .}
\expandafter\gdef\csname CGTip@prop:rounding\endcsname{Proposition 5.5 (Rounding the corrected count misassigns a positive proportion) (p. 18): There are absolute constants c\_*>0, kappa\_0>0 and T\_0 such that for T>= T\_0 and 2<= X<= T\textasciicircum{}c\_*, [ (1)/(T)meastin[T,2T]:floor(F\_X(t)+(1)/(2))!= N(t)>=kappa\_0 . ]}
\expandafter\gdef\csname CGTip@rem:comparison_LeClair\endcsname{Remark 5.6 (p. 18): Since logno.p<= X=log X+O(log log X), the condition of Corollary 5.4 says that the logarithm of the number of primes must grow proportionally to log T. A fixed number of primes, as in the computations reported in [ref, S XXI, Note added 7], leaves almost all zeros violating the corrected equation by more than any fixe...}
\expandafter\gdef\csname CGTip@rem:localization\endcsname{Remark 5.7 (Localization of the ordinates) (p. 19): If M\_X<=(1)/(2)L(T), then F\_X'>=L-M\_X>=(1)/(2)L(T) on [T,2T], so F\_X is strictly increasing there. For kin I\_T with F\_X(T)<= k-(1)/(2)<= F\_X(2T) the equation F\_X(t)=k-(1)/(2) has exactly one root t\_k\textasciicircum{}Xin[T,2T] (intermediate value theorem), and u\_k(X)=F\_X(t\_k\textasciicircum{}X)-F\_X(gamma\_k) gives |gamma\_k-t\_k\textasciicircum{}X|<=2|u\_k(X)|/L(T) (mea...}
\expandafter\gdef\csname CGTip@subsec:main_lemmas\endcsname{Section 5.1: Lemmas on prime polynomials (p. 19)}
\expandafter\gdef\csname CGTip@eq:D_def\endcsname{Eq. (5.4), p. 19: W\_1(t)=-(1)/(pi)sum\_pin J(sin(tlog p))/(sqrt(p)), V=sum\_pin J(1)/(p),; W=S\textasciicircum{}rm pr\_Y\_0-S\_X=+- W\_1-S\_X\textasciicircum{}(2)-S\_X\textasciicircum{}(>=3),}
\expandafter\gdef\csname CGTip@lem:V_size\endcsname{Lemma 5.8 (p. 19): V=log(1/c)+O(1) uniformly for log 2/log T<= c<=(1)/(2).}
\expandafter\gdef\csname CGTip@lem:trig_MV\endcsname{Lemma 5.9 (p. 19): For real z\_n, 2<= n<= Y, int\_T\textasciicircum{}2T(sum\_n<= Yz\_nsin(tlog n))\textasciicircum{}2dt=(T)/(2)sum\_n<= Yz\_n\textasciicircum{}2+O(sum\_n<= Ynz\_n\textasciicircum{}2+(sum\_n<= Y|z\_n|)\textasciicircum{}2), and the same with cos in place of sin.}
\expandafter\gdef\csname CGTip@eq:trig_MV\endcsname{Eq. (5.5), p. 19: int\_T\textasciicircum{}2T(sum\_n<= Yz\_nsin(tlog n))\textasciicircum{}2dt=(T)/(2)sum\_n<= Yz\_n\textasciicircum{}2+O(sum\_n<= Ynz\_n\textasciicircum{}2+(sum\_n<= Y|z\_n|)\textasciicircum{}2),}
\expandafter\gdef\csname CGTip@lem:D_moments\endcsname{Lemma 5.10 (p. 19): With absolute implied constants, int\_T\textasciicircum{}2T W\textasciicircum{}2 dt =(T)/(2pi\textasciicircum{}2)V+O(T(1+sqrt(V))), ; int\_T\textasciicircum{}2T W\textasciicircum{}4 dt << T(1+V)\textasciicircum{}2, ; int\_T\textasciicircum{}2TS\_X'(t)\textasciicircum{}2 dt << T(1+log X)\textasciicircum{}2 .}
\expandafter\gdef\csname CGTip@eq:D_second\endcsname{Eq. (5.6), p. 19: int\_T\textasciicircum{}2T W\textasciicircum{}2 dt =(T)/(2pi\textasciicircum{}2)V+O(T(1+sqrt(V))),}
\expandafter\gdef\csname CGTip@eq:D_fourth\endcsname{Eq. (5.7), p. 19: int\_T\textasciicircum{}2T W\textasciicircum{}4 dt << T(1+V)\textasciicircum{}2,}
\expandafter\gdef\csname CGTip@eq:SX_derivative\endcsname{Eq. (5.8), p. 19: int\_T\textasciicircum{}2TS\_X'(t)\textasciicircum{}2 dt << T(1+log X)\textasciicircum{}2 .}
\expandafter\gdef\csname CGTip@subsec:main_proof\endcsname{Section 5.2: Proof of Theorem 5.1 (p. 20)}
\expandafter\gdef\csname CGTip@eq:E0_bounds\endcsname{Eq. (5.9), p. 20: int\_T\textasciicircum{}2TE\_0\textasciicircum{}4 dt<< T,int\_T\textasciicircum{}2TE\_0\textasciicircum{}2 dt<= T\textasciicircum{}1/2(int\_T\textasciicircum{}2TE\_0\textasciicircum{}4 dt)\textasciicircum{}1/2<< T .}
\expandafter\gdef\csname CGTip@eq:J_asymptotic\endcsname{Eq. (5.10), p. 20: E\_2=(T)/(2pi\textasciicircum{}2)V+O(T(1+sqrt(V))),}
\expandafter\gdef\csname CGTip@eq:R_fourth\endcsname{Eq. (5.11), p. 20: int\_T\textasciicircum{}2TE\_X\textasciicircum{}4 dt<< T(1+V)\textasciicircum{}2 .}
\expandafter\gdef\csname CGTip@eq:main_step1\endcsname{Eq. (5.12), p. 20: sum\_kin I\_Tu\_k(X)\textasciicircum{}2=int\_T\textasciicircum{}2TE\_X\textasciicircum{}2L dt+int\_T\textasciicircum{}2TE\_X\textasciicircum{}2S\_X' dt-(no.I\_T)/(12)+O(T\textasciicircum{}3/4).}
\expandafter\gdef\csname CGTip@eq:main_step2\endcsname{Eq. (5.13), p. 21: |int\_T\textasciicircum{}2TE\_X\textasciicircum{}2S\_X' dt|<=(int\_T\textasciicircum{}2TE\_X\textasciicircum{}4)\textasciicircum{}1/2(int\_T\textasciicircum{}2TS\_X'\textasciicircum{}2)\textasciicircum{}1/2<< T(1+V)(1+log X).}
\expandafter\gdef\csname CGTip@eq:main_step3\endcsname{Eq. (5.14), p. 21: (1)/(no.I\_T)sum\_kin I\_Tu\_k(X)\textasciicircum{}2=(V)/(2pi\textasciicircum{}2)+O(1+sqrt(V)+(1+V)/(log T)+(1+V) c),}
\expandafter\gdef\csname CGTip@subsec:variable_cutoff\endcsname{Section 5.3: A cutoff varying with the index (p. 21)}
\expandafter\gdef\csname CGTip@lem:maximal_primes\endcsname{Lemma 5.11 (p. 21): Let T>= T\_0 and 2<= X\_-<X\_+<=(T/3)\textasciicircum{}1/2 with X\_+<=4X\_-; let V\_+-=sum\_X\_-<p<= X\_+p\textasciicircum{}-1 and [ Q(t)=max\_X\_-<= y<= X\_+|sum\_X\_-<p<= yp\textasciicircum{}-1/2-it| . ] Then Q is Lipschitz, | Q'(t)|<=max\_X\_-<= y<= X\_+|sum\_X\_-<p<= ylog p p\textasciicircum{}-1/2-it| for almost every t, and int\_T-3\textasciicircum{}2T+3 Q\textasciicircum{}4 dt<< TV\_+-\textasciicircum{}2,int\_T-3\textasciicircum{}2T+3| Q'|\textasciicircum{}4 dt<< T(log X\_+)\textasciicircum{}4V\_+-\textasciicircum{}2...}
\expandafter\gdef\csname CGTip@eq:maximal_L4\endcsname{Eq. (5.15), p. 21: int\_T-3\textasciicircum{}2T+3 Q\textasciicircum{}4 dt<< TV\_+-\textasciicircum{}2,int\_T-3\textasciicircum{}2T+3| Q'|\textasciicircum{}4 dt<< T(log X\_+)\textasciicircum{}4V\_+-\textasciicircum{}2,}
\expandafter\gdef\csname CGTip@eq:maximal_points\endcsname{Eq. (5.16), p. 21: < Q(gamma\_k)\textasciicircum{}2>\_I\_T<< V\_+-(1+(log X\_+)/(log T)sqrt(log log T))+(X\_+)/(T),< Q(tau\_k)\textasciicircum{}2>\_ I\textasciicircum{}tau\_T<< V\_+- .}
\expandafter\gdef\csname CGTip@eq:maximal_stieltjes\endcsname{Eq. (5.17), p. 22: sum\_kin I\_T Q(gamma\_k)\textasciicircum{}2=int\_T\textasciicircum{}2T Q\textasciicircum{}2L dt+[ Q\textasciicircum{}2G]\_T\textasciicircum{}2T-2int\_T\textasciicircum{}2TG Q Q' dt,}
\expandafter\gdef\csname CGTip@thm:variable_cutoff\endcsname{Theorem 5.12 (p. 22): For (log 2)/(log T)<= c<=(1)/(2) let X\_k=max 2,(tau\_k/2pi)\textasciicircum{}c and tilde u\_k=s\_k-S\_X\_k(gamma\_k). Uniformly in c, |<tilde u\_k\textasciicircum{}2>\_I\_T-(1)/(2pi\textasciicircum{}2)log(1)/(c)|<<1+sqrt(log(1/c)) . For fixed c, <(tilde u\_k-u\_k(X\_-))\textasciicircum{}2>\_I\_T->0, with X\_- as in the proof.}
\expandafter\gdef\csname CGTip@eq:variable_cutoff\endcsname{Eq. (5.18), p. 22: |<tilde u\_k\textasciicircum{}2>\_I\_T-(1)/(2pi\textasciicircum{}2)log(1)/(c)|<<1+sqrt(log(1/c)) .}
\expandafter\gdef\csname CGTip@sec:accounting\endcsname{Section 6: Theorem C: exact accounting for the one-step clock (p. 23)}
\expandafter\gdef\csname CGTip@eq:two_responses\endcsname{Eq. (6.1), p. 23: R\textasciicircum{}gamma\_n(I)=(2)/(a\_n)< s\_ksin(gamma\_klog n)>\_I, R\textasciicircum{}tau\_n(I)=(2)/(a\_n)< s\_ksin(tau\_klog n)>\_I,}
\expandafter\gdef\csname CGTip@eq:Cr_def\endcsname{Eq. (6.2), p. 23: Cr\_A(I)=< S\_A(tau\_k)\textasciicircum{}2>\_I-v\_A .}
\expandafter\gdef\csname CGTip@thm:accounting\endcsname{Theorem 6.1 (Finite-window accounting) (p. 23): Let I be a finite nonempty set of indices and A a set of prime powers with sum\_nin A|a\_n|<inf. item For every k and X>=2, r\_k(X)=u\_k(X)+(S\_X(gamma\_k)-S\_X(tau\_k))+chi\_k,; chi\_k=(gamma\_k)/(2pi)((tau\_k)/(gamma\_k)-1-log(tau\_k)/(gamma\_k)),(delta\_k\textasciicircum{}2)/(4pi max(gamma\_k,tau\_k))<=chi\_k<=(delta\_k\textasciicircum{}2)/(4pi min(gamma\_k,tau\_k)),;...}
\expandafter\gdef\csname CGTip@eq:one_step_decomposition\endcsname{Eq. (6.3), p. 23: r\_k(X)=u\_k(X)+(S\_X(gamma\_k)-S\_X(tau\_k))+chi\_k,; chi\_k=(gamma\_k)/(2pi)((tau\_k)/(gamma\_k)-1-log(tau\_k)/(gamma\_k)),(delta\_k\textasciicircum{}2)/(4pi max(gamma\_k,tau\_k))<=chi\_k<=(delta\_k\textasciicircum{}2)/(4pi min(gamma\_k,tau\_k)),; |S\_X(gamma\_k)-S\_X(tau\_k)|<= M\_X|delta\_k|.}
\expandafter\gdef\csname CGTip@eq:exact_gain\endcsname{Eq. (6.4), p. 23: < s\_k\textasciicircum{}2-(s\_k-lambda S\_A(tau\_k))\textasciicircum{}2>\_I =sum\_nin Aa\_n\textasciicircum{}2(lambda R\textasciicircum{}tau\_n(I)-(lambda\textasciicircum{}2)/(2))-lambda\textasciicircum{}2Cr\_A(I).}
\expandafter\gdef\csname CGTip@eq:full_factor_gain\endcsname{Eq. (6.5), p. 23: < s\_k\textasciicircum{}2-(s\_k-S\_X(tau\_k))\textasciicircum{}2>\_I=D\_X(I)-Cr\_X(I),; D\_X(I)=sum\_nin A\_Xa\_n\textasciicircum{}2(R\textasciicircum{}tau\_n(I)-(1)/(2))=2< s\_kS\_X(tau\_k)>\_I-v\_X,}
\expandafter\gdef\csname CGTip@eq:physical_gain\endcsname{Eq. (6.6), p. 24: <(L(tau\_k)delta\_k)\textasciicircum{}2-r\_k(X)\textasciicircum{}2>\_I=D\_X(I)-Cr\_X(I)+2<chi\_kS\_X(tau\_k)>\_I,; |<chi\_kS\_X(tau\_k)>\_I|<=sigma\_Xmax\_kin Ichi\_k .}
\expandafter\gdef\csname CGTip@rem:accounting_meaning\endcsname{Remark 6.2 (p. 24): (6.5) is an exact sign criterion for the gain of the full Euler phase at the clock points: the gain is positive exactly when the diagonal responses outweigh the cross term. No independence of the phases tau\_klog p is assumed. A threshold R\textasciicircum{}tau\_p=(1)/(2) describes a single diagonal prime contribution and does not det...}
\expandafter\gdef\csname CGTip@rem:matrix_projection\endcsname{Remark 6.3 (Optimal trigonometric correction) (p. 24): For a finite set A of frequencies let varsigma\_n(k)=sin(tau\_klog n), Q=(<varsigma\_n\_1varsigma\_n\_2>\_I)\_n\_1,n\_2in A, b=(< s\_kvarsigma\_n(k)>\_I)\_nin A and Upsilon\_lambda=sum\_nlambda\_nvarsigma\_n. Ordinary least squares gives <(s\_k-Upsilon\_lambda(k))\textasciicircum{}2>\_I=< s\_k\textasciicircum{}2>\_I-2lambda\textasciicircum{} T b+lambda\textasciicircum{} T Qlambda, minimized by lambda= Q\textasciicircum{}d...}
\expandafter\gdef\csname CGTip@rem:transport\endcsname{Remark 6.4 (Clock point against zero) (p. 24): For a prime power n, sin(gamma\_kl\_n)-sin(tau\_kl\_n)=int\_0\textasciicircum{}1(d)/(dy)sin((tau\_k-ydelta\_k)l\_n) dy, whence, exchanging the finite mean and the integral, R\textasciicircum{}gamma\_n(I)-R\textasciicircum{}tau\_n(I)=-(2l\_n)/(a\_n)int\_0\textasciicircum{}1< s\_kdelta\_kcos ((tau\_k-ydelta\_k)l\_n)>\_I dy . For fixed n and I=I\_T the right side is O(<|s\_kdelta\_k|>\_I\_T)=O(log log T/log T...}
\expandafter\gdef\csname CGTip@eq:response_transport\endcsname{Eq. (6.7), p. 24: R\textasciicircum{}gamma\_n(I)-R\textasciicircum{}tau\_n(I)=-(2l\_n)/(a\_n)int\_0\textasciicircum{}1< s\_kdelta\_kcos ((tau\_k-ydelta\_k)l\_n)>\_I dy .}
\expandafter\gdef\csname CGTip@rem:fixed_sup\endcsname{Remark 6.5 (Fixed prime sets) (p. 24): For fixed X, |delta\_k\textasciicircum{}X-delta\_k|<=sigma\_X/L(tau\_k)->0 and |u\_k(X)-s\_k|<=sigma\_X. A fixed set of primes changes neither limsup\_k|delta\_k| nor the unboundedness of the defect.}
\expandafter\gdef\csname CGTip@sec:computable\endcsname{Section 7: Theorem D: a computable clock that approaches the zeros (p. 25)}
\expandafter\gdef\csname CGTip@eq:r_decomposition\endcsname{Eq. (7.1), p. 25: r\_k(X)=u\_k(X)+ t\_k(X)+chi\_k, t\_k(X)=S\_X(gamma\_k)-S\_X(tau\_k),}
\expandafter\gdef\csname CGTip@lem:SX_fourth\endcsname{Lemma 7.1 (p. 25): For T>= T\_0 and 2<= X<= T\textasciicircum{}1/2, [ int\_T-3\textasciicircum{}2T+3S\_X'(t)\textasciicircum{}4 dt<< T(1+log X)\textasciicircum{}4,int\_T-3\textasciicircum{}2T+3S\_X''(t)\textasciicircum{}4 dt<< T(1+log X)\textasciicircum{}8 . ]}
\expandafter\gdef\csname CGTip@lem:transport\endcsname{Lemma 7.2 (Transport) (p. 25): For T>= T\_0 and 2<= X<= T\textasciicircum{}1/2, < t\_k(X)\textasciicircum{}2>\_I\_T<<(log log T (1+log X)\textasciicircum{}2)/(log\textasciicircum{}2T).}
\expandafter\gdef\csname CGTip@eq:transport\endcsname{Eq. (7.2), p. 25: < t\_k(X)\textasciicircum{}2>\_I\_T<<(log log T (1+log X)\textasciicircum{}2)/(log\textasciicircum{}2T).}
\expandafter\gdef\csname CGTip@eq:delta_fourth\endcsname{Eq. (7.3), p. 25: <delta\_k\textasciicircum{}4>\_I\_T=< s\_k\textasciicircum{}4L(xi\_k)\textasciicircum{}-4>\_I\_T<<((log log T)\textasciicircum{}2)/(log\textasciicircum{}4T).}
\expandafter\gdef\csname CGTip@eq:transport_maximal\endcsname{Eq. (7.4), p. 25: | t\_k(X)|=|int\_tau\_k\textasciicircum{}gamma\_k h(y) dy|<=int\_tau\_k-|delta\_k|\textasciicircum{}tau\_k+|delta\_k|| h|<=2|delta\_k| M h(tau\_k)}
\expandafter\gdef\csname CGTip@thm:computable_clock\endcsname{Theorem 7.3 (Theorem D) (p. 26): There are absolute constants T\_0 and C such that for T>= T\_0 and 2<= X<= T\textasciicircum{}1/2, with c=log X/log T, |< r\_k(X)\textasciicircum{}2>\_I\_T-(1)/(2pi\textasciicircum{}2)log(1)/(c)|<= C(1+sqrt(log(1/c)))(1+c\textasciicircum{}2log log T).}
\expandafter\gdef\csname CGTip@eq:computable_clock\endcsname{Eq. (7.5), p. 26: |< r\_k(X)\textasciicircum{}2>\_I\_T-(1)/(2pi\textasciicircum{}2)log(1)/(c)|<= C(1+sqrt(log(1/c)))(1+c\textasciicircum{}2log log T).}
\expandafter\gdef\csname CGTip@cor:computable_clock\endcsname{Corollary 7.4 (p. 26): Let X=T\textasciicircum{}1/sqrt(log log T). Then < r\_k(X)\textasciicircum{}2>\_I\_T =(1)/(4pi\textasciicircum{}2)log log log T+O(sqrt(log log log T)), ; (1)/(no.I\_T)sum\_kin I\_T(tau\_k\textasciicircum{}X-gamma\_k)\textasciicircum{}2 =(log log log T+O(sqrt(log log log T)))/(log\textasciicircum{}2T), while (1)/(no.I\_T)sum\_kin I\_T(tau\_k-gamma\_k)\textasciicircum{}2=(2+o(1))log log T/log\textasciicircum{}2T; the ratio is asymptotic to log log log T/(2log log T).}
\expandafter\gdef\csname CGTip@eq:computable_units\endcsname{Eq. (7.6), p. 26: < r\_k(X)\textasciicircum{}2>\_I\_T =(1)/(4pi\textasciicircum{}2)log log log T+O(sqrt(log log log T)),}
\expandafter\gdef\csname CGTip@eq:computable_height\endcsname{Eq. (7.7), p. 26: (1)/(no.I\_T)sum\_kin I\_T(tau\_k\textasciicircum{}X-gamma\_k)\textasciicircum{}2 =(log log log T+O(sqrt(log log log T)))/(log\textasciicircum{}2T),}
\expandafter\gdef\csname CGTip@cor:local_clock\endcsname{Corollary 7.5 (A cutoff read off the clock) (p. 26): For tau\_k>e\textasciicircum{}e let X\_k\textasciicircum{}rm loc=tau\_k\textasciicircum{}1/sqrt(log log tau\_k) and tau\_k\textasciicircum{}rm loc=tau\_k-S\_X\_k\textasciicircum{}rm loc(tau\_k)/L(tau\_k). Then [ (1)/(no. I\_T\textasciicircum{}tau)sum\_kin I\_T\textasciicircum{}tau(tau\_k\textasciicircum{}rm loc-gamma\_k)\textasciicircum{}2=(log log log T+O(sqrt(log log log T)))/(log\textasciicircum{}2T). ]}
\expandafter\gdef\csname CGTip@prop:monotone\endcsname{Proposition 7.6 (The root of the corrected equation in the monotone regime) (p. 27): Let X=X(T)>=2 with M\_X<=(1)/(2)L(T), and let I'\_T be the set of kin I\_T with F\_X(T)<= k-(1)/(2)<= F\_X(2T), so that the root t\_k\textasciicircum{}Xin[T,2T] of Remark 5.7 exists. Then X<<log\textasciicircum{}2T, no.(I\_T I'\_T)<<log T, and [ (1)/(no.I\_T)sum\_kin I'\_T(t\_k\textasciicircum{}X-gamma\_k)\textasciicircum{}2>=(4)/(9L(2T)\textasciicircum{}2)((1)/(2pi\textasciicircum{}2)log log T-O(sqrt(log log T)))>>(log log T)/(log\textasciicircum{}...}
\expandafter\gdef\csname CGTip@rem:newton\endcsname{Remark 7.7 (The Newton term and the restriction on c) (p. 27): Let xi\textasciicircum{}X\_k be a point between gamma\_k and tau\_k with t\_k(X)=-S\_X'(xi\_k\textasciicircum{}X)delta\_k (mean value theorem). Since delta\_k=(r\_k(X)+S\_X(tau\_k))/L(tau\_k) by (2.12), (7.1) becomes the exact relation r\_k(X)(1+(S\_X'(xi\_k\textasciicircum{}X))/(L(tau\_k)))=u\_k(X)+chi\_k-y\_k S\_X'(xi\_k\textasciicircum{}X). The term y\_kS\_X'(xi\_k\textasciicircum{}X) is the second-order term of Newton'...}
\expandafter\gdef\csname CGTip@eq:newton_exact\endcsname{Eq. (7.8), p. 27: r\_k(X)(1+(S\_X'(xi\_k\textasciicircum{}X))/(L(tau\_k)))=u\_k(X)+chi\_k-y\_k S\_X'(xi\_k\textasciicircum{}X).}
\expandafter\gdef\csname CGTip@sec:iteration\endcsname{Section 8: Iterating the Newton step (p. 28)}
\expandafter\gdef\csname CGTip@subsec:recursion\endcsname{Section 8.1: The exact recursion (p. 28)}
\expandafter\gdef\csname CGTip@eq:newton_iteration\endcsname{Eq. (8.1), p. 28: x\textasciicircum{}(j)\_k=x\textasciicircum{}(j-1)\_k-(N\_0(x\textasciicircum{}(j-1)\_k)+S\_X\_j(x\textasciicircum{}(j-1)\_k)-(k-(1)/(2)))/(L(x\textasciicircum{}(j-1)\_k)), e\_j,k=x\textasciicircum{}(j)\_k-gamma\_k,}
\expandafter\gdef\csname CGTip@prop:recursion\endcsname{Proposition 8.1 (Exact recursion) (p. 28): Let j>=1 and suppose x\textasciicircum{}(j-1)\_k>2pi. Then L(x\textasciicircum{}(j-1)\_k) e\_j,k=u\_k(X\_j)+(S\_X\_j(gamma\_k)-S\_X\_j(x\textasciicircum{}(j-1)\_k))+chi\_j,k, 0<=chi\_j,k<=(e\_j-1,k\textasciicircum{}2)/(4pi min(gamma\_k,x\textasciicircum{}(j-1)\_k)) . Let varphi\_jin C\textasciicircum{}inf(R), 0<=varphi\_j<=1, be 1 on an interval containing all segments between gamma\_k and x\textasciicircum{}(j-1)\_k, kin I, and g\_j= M(varphi\_jS'\_X\_j)....}
\expandafter\gdef\csname CGTip@eq:recursion_identity\endcsname{Eq. (8.2), p. 28: L(x\textasciicircum{}(j-1)\_k) e\_j,k=u\_k(X\_j)+(S\_X\_j(gamma\_k)-S\_X\_j(x\textasciicircum{}(j-1)\_k))+chi\_j,k, 0<=chi\_j,k<=(e\_j-1,k\textasciicircum{}2)/(4pi min(gamma\_k,x\textasciicircum{}(j-1)\_k)) .}
\expandafter\gdef\csname CGTip@eq:recursion_norms\endcsname{Eq. (8.3), p. 28: ||L(x\textasciicircum{}(j-1))e\_j||\_2r,I<=||u(X\_j)||\_2r,I+2 ||e\_j-1||\_4r,I ||g\_j(x\textasciicircum{}(j-1))||\_4r,I+||chi\_j||\_2r,I,; ||chi\_j||\_2r,I<=(||e\_j-1||\_4r,I\textasciicircum{}2)/(4pi min\_kin Imin(gamma\_k,x\textasciicircum{}(j-1)\_k)) .}
\expandafter\gdef\csname CGTip@lem:higher_moments\endcsname{Lemma 8.2 (p. 29): There is an absolute constant C such that for every integer r>=1, every T with T\textasciicircum{}1/(160r)>=2 and every 2<= X<= T\textasciicircum{}1/(160r), with V as in Lemma 5.8 computed with Y\_2r=T\textasciicircum{}1/(160r) in place of Y\_0, [ int\_T\textasciicircum{}2TE\_X\textasciicircum{}4r dt<=(C(r\textasciicircum{}2+sqrt(r(1+V))))\textasciicircum{}4rT, < u\_k(X)\textasciicircum{}2r>\_I\_T<=(C(r\textasciicircum{}2+sqrt(r(1+V))))\textasciicircum{}2r, ] and for every r>=1 and 2<= X<=...}
\expandafter\gdef\csname CGTip@lem:multiplicity\endcsname{Lemma 8.3 (Multiplicity) (p. 29): int\_R tilde n(u)\textasciicircum{}2 du=sum\_j,kin I\_T(b\_T-|rho\_j-rho\_k|)\_+<< b\_T no.I\_T .}
\expandafter\gdef\csname CGTip@eq:close_pair\endcsname{Eq. (8.4), p. 30: |y\_k+m-y\_k|>=(m-1)b .}
\expandafter\gdef\csname CGTip@lem:sampling_rho\endcsname{Lemma 8.4 (Sampling at points of bounded multiplicity) (p. 30): Let varrho\_k, kin I\_T, be real points of [T-2T\textasciicircum{}1/2,2T+2T\textasciicircum{}1/2] with multiplicity function tilde n(u)=no.kin I\_T:u- b\_T<=varrho\_k<= u satisfying int\_R tilde n\textasciicircum{}2<=kappa b\_Tno.I\_T for some kappa >=1. Let r>=1 be an integer, 2<= X<= T\textasciicircum{}1/(4r), varphi a cutoff equal to 1 on [T-2T\textasciicircum{}1/2,2T+2T\textasciicircum{}1/2] as in Lemma 7.2, and g= M(varphi...}
\expandafter\gdef\csname CGTip@thm:second_step\endcsname{Theorem 8.5 (Second Newton step) (p. 31): Let H\_0=log log T, H\_1=log H\_0, H\_2=log H\_1, X\_1=T\textasciicircum{}H\_0\textasciicircum{}-1/2, X\_2=T\textasciicircum{}H\_1\textasciicircum{}-1/2, and let x\textasciicircum{}(2)\_k be given by (8.1). Then <L(rho\_k)\textasciicircum{}2e\_2,k\textasciicircum{}2>\_I\_T=(1)/(4pi\textasciicircum{}2)H\_2+O(1+sqrt(H\_2)),; (1)/(no.I\_T)sum\_kin I\_T(x\textasciicircum{}(2)\_k-gamma\_k)\textasciicircum{}2=(log log log log T+O(sqrt(log log log log T)))/(log\textasciicircum{}2T).}
\expandafter\gdef\csname CGTip@eq:second_step\endcsname{Eq. (8.5), p. 31: <L(rho\_k)\textasciicircum{}2e\_2,k\textasciicircum{}2>\_I\_T=(1)/(4pi\textasciicircum{}2)H\_2+O(1+sqrt(H\_2)),; (1)/(no.I\_T)sum\_kin I\_T(x\textasciicircum{}(2)\_k-gamma\_k)\textasciicircum{}2=(log log log log T+O(sqrt(log log log log T)))/(log\textasciicircum{}2T).}
\expandafter\gdef\csname CGTip@eq:e1_fourth\endcsname{Eq. (8.6), p. 31: < e\_1,k\textasciicircum{}4>\_I\_T<<(H\_1\textasciicircum{}2)/(log\textasciicircum{}4T).}
\expandafter\gdef\csname CGTip@subsec:further_steps\endcsname{Section 8.5: All further steps: the exact identity, the local deformation, and the multiplicity (p. 32)}
\expandafter\gdef\csname CGTip@lem:second_identity\endcsname{Lemma 8.6 (The second-level identity) (p. 32): For j>=1, with Y\textasciicircum{}(j)\_k=tau\_k-x\textasciicircum{}(j)\_k and chi''\_j,k=int\_x\textasciicircum{}(j-1)\_k\textasciicircum{}tau\_k(L(tau\_k)-L(y)) dy, Y\textasciicircum{}(j)\_k=(1-(L(tau\_k))/(L(x\textasciicircum{}(j-1)\_k)))Y\textasciicircum{}(j-1)\_k+(S\_X\_j(x\textasciicircum{}(j-1)\_k)+chi''\_j,k)/(L(x\textasciicircum{}(j-1)\_k)),; henceL(tau\_k)Y\textasciicircum{}(j)\_k=S\_X\_j(x\textasciicircum{}(j-1)\_k)+R\_j,k, with |R\_j,k|<< T\textasciicircum{}c\_j-1 if |Y\textasciicircum{}(j-1)\_k|<= T\textasciicircum{}c\_j-1/2 and c\_j-1<= c\_j (true by (2.11) and ind...}
\expandafter\gdef\csname CGTip@eq:Y_identity\endcsname{Eq. (8.7), p. 32: Y\textasciicircum{}(j)\_k=(1-(L(tau\_k))/(L(x\textasciicircum{}(j-1)\_k)))Y\textasciicircum{}(j-1)\_k+(S\_X\_j(x\textasciicircum{}(j-1)\_k)+chi''\_j,k)/(L(x\textasciicircum{}(j-1)\_k)),; henceL(tau\_k)Y\textasciicircum{}(j)\_k=S\_X\_j(x\textasciicircum{}(j-1)\_k)+R\_j,k,}
\expandafter\gdef\csname CGTip@prop:deformation\endcsname{Proposition 8.7 (Local deformation of the second clock) (p. 32): Let X\_1=T\textasciicircum{}c\_1, X\_2=T\textasciicircum{}c\_2 with c\_1=H\_0\textasciicircum{}-1/2, c\_2=H\_1\textasciicircum{}-1/2, and d\_j,k=x\textasciicircum{}(j)\_k+1-x\textasciicircum{}(j)\_k (d\_0,k=tau\_k+1-tau\_k). Then, for T>= T\_0, <|(d\_2,k-d\_0,k)/(d\_0,k)|\textasciicircum{}4>\_I\_T<< c\_2\textasciicircum{}4=(1)/(H\_1\textasciicircum{}2), (1)/(no.I\_T)no.kin I\_T:|(d\_2,k)/(d\_0,k)-1|>alpha <<(1)/(alpha\textasciicircum{}4H\_1\textasciicircum{}2)(0<alpha <=1), in particular no.kin I\_T:d\_2,k<=0<<no.I\_T/(log log log T)\textasciicircum{}2:...}
\expandafter\gdef\csname CGTip@eq:deformation\endcsname{Eq. (8.8), p. 32: <|(d\_2,k-d\_0,k)/(d\_0,k)|\textasciicircum{}4>\_I\_T<< c\_2\textasciicircum{}4=(1)/(H\_1\textasciicircum{}2), (1)/(no.I\_T)no.kin I\_T:|(d\_2,k)/(d\_0,k)-1|>alpha <<(1)/(alpha\textasciicircum{}4H\_1\textasciicircum{}2)(0<alpha <=1),}
\expandafter\gdef\csname CGTip@lem:zero_pairs\endcsname{Lemma 8.8 (Near pairs and triples of ordinates) (p. 33): Let b= b\_T, M=no.I\_T and, for Lambda >=1, n\_Lambda(u)=no.l:|gamma\_l-u|<=Lambda b. There are absolute constants C\_1,C\_2 such that for T>= T\_0 and every Lambda >=1, no.(k,l)in I\_T\textasciicircum{}2:|gamma\_k-gamma\_l|<=Lambda b<= C\_1Lambda M,sum\_kin I\_Tn\_Lambda(gamma\_k)\textasciicircum{}2<= C\_2Lambda\textasciicircum{}2M .}
\expandafter\gdef\csname CGTip@eq:zero_pairs\endcsname{Eq. (8.9), p. 33: no.(k,l)in I\_T\textasciicircum{}2:|gamma\_k-gamma\_l|<=Lambda b<= C\_1Lambda M,sum\_kin I\_Tn\_Lambda(gamma\_k)\textasciicircum{}2<= C\_2Lambda\textasciicircum{}2M .}
\expandafter\gdef\csname CGTip@thm:multiplicity_all\endcsname{Theorem 8.9 (Multiplicity of the corrected clock points) (p. 33): Let j>=0 and let x\textasciicircum{}(j)\_k be given by (8.1) with cutoffs X\_i=T\textasciicircum{}c\_i, c\_i<=(1)/(2), so that x\textasciicircum{}(j)\_kin[T-2T\textasciicircum{}1/2,2T+2T\textasciicircum{}1/2] for kin I\_T and T>= T\_0 (Lemma 8.6). Let tilde n\_j(u)=no.kin I\_T:u- b\_T<= x\textasciicircum{}(j)\_k<= u, and let A>=1, H>=1 be such that ||L(tau\_k)e\_j,k||\_4,I\_T<= Asqrt(H). Then, with an absolute constant C\_3, int\_R ti...}
\expandafter\gdef\csname CGTip@eq:multiplicity_all\endcsname{Eq. (8.10), p. 33: int\_R tilde n\_j(u)\textasciicircum{}2 du<= C\_3 Asqrt(H) b\_T no.I\_T .}
\expandafter\gdef\csname CGTip@eq:cutoff_law\endcsname{Eq. (8.11), p. 34: H\_j=2log(1)/(c\_j)(j>=1), c\_j=C\_* H\_j-1\textasciicircum{}-9/16(j>=3),}
\expandafter\gdef\csname CGTip@thm:uniform_iteration\endcsname{Theorem 8.10 (All Newton steps, uniformly) (p. 34): There are absolute constants C\_*, H\textasciicircum{}*, C and T\_0 such that, with the cutoffs (8.11), for every T>= T\_0 and every 1<= m<= m\textasciicircum{}*, <L(tau\_k)\textasciicircum{}2e\_m,k\textasciicircum{}2>\_I\_T=(H\_m)/(4pi\textasciicircum{}2)+vartheta\_m C(1+sqrt(H\_m)),|vartheta\_m|<=1, and ||L(tau\_k)e\_j,k||\_4,I\_T<= Csqrt(H\_j) for 1<= j<m\textasciicircum{}*. In particular, at m=m\textasciicircum{}*=log\textasciicircum{}*T+O(1), (1)/(no.I\_T)sum\_kin...}
\expandafter\gdef\csname CGTip@eq:uniform_iteration\endcsname{Eq. (8.12), p. 34: <L(tau\_k)\textasciicircum{}2e\_m,k\textasciicircum{}2>\_I\_T=(H\_m)/(4pi\textasciicircum{}2)+vartheta\_m C(1+sqrt(H\_m)),|vartheta\_m|<=1,}
\expandafter\gdef\csname CGTip@eq:final_clock\endcsname{Eq. (8.13), p. 34: (1)/(no.I\_T)sum\_kin I\_T(x\textasciicircum{}(m\textasciicircum{}*)\_k-gamma\_k)\textasciicircum{}2<=(C')/(log\textasciicircum{}2T):}
\expandafter\gdef\csname CGTip@eq:uniform_recursion\endcsname{Eq. (8.14), p. 35: A\_j(r)<= C(sqrt(r)+(r\textasciicircum{}2)/(sqrt(H\_j)))+(C\_5sqrt(r))/(sqrt(H\_j))max (A\_j-1(2r),A\_j-1(2r)\textasciicircum{}9/8),}
\expandafter\gdef\csname CGTip@eq:induction_claim\endcsname{Eq. (8.15), p. 35: A\_j(r)<= Br\textasciicircum{}2(0<= j<= m-1, 1<= r<=2\textasciicircum{}m-j).}
\expandafter\gdef\csname CGTip@cor:third_step\endcsname{Corollary 8.11 (Third step) (p. 36): With c\_3=C\_* H\_2\textasciicircum{}-9/16 and H\_3=(9)/(8)log H\_2+2log(1/C\_*), [ <L(tau\_k)\textasciicircum{}2e\_3,k\textasciicircum{}2>\_I\_T=(H\_3)/(4pi\textasciicircum{}2)+O(1+sqrt(H\_3)),(1)/(no.I\_T)sum\_kin I\_T(x\textasciicircum{}(3)\_k-gamma\_k)\textasciicircum{}2=(H\_3+O(sqrt(H\_3)))/(log\textasciicircum{}2T), ] and the multiplicity of the second clock satisfies int tilde n\_2\textasciicircum{}2llsqrt(H\_2) b\_Tno.I\_T.}
\expandafter\gdef\csname CGTip@tab:iteration\endcsname{Table 2 (p. 36): Iterated Newton steps over 2\textasciicircum{}19<gamma\_k<=2\textasciicircum{}20 with fixed cutoff exponents (c\_1,c\_2,c\_3): mean squares of L e\_j in units of zeros (j=0 is the bare clock), fourth moment and root mean square of the relative deformation (d\_2,k-d\_0,k)/d\_0,k, proportion of inversions d\_2,k<=0, and the multiplicity integrals of x\textasciicircum{}(1) and...}
\expandafter\gdef\csname CGTip@subsec:twist\endcsname{Section 8.6: The transport identity in a prime frequency (p. 36)}
\expandafter\gdef\csname CGTip@eq:twist_def\endcsname{Eq. (8.16), p. 36: I\_T(w)=(1)/(T)int\_T\textasciicircum{}2Texp (iwt+(iw G(t))/(L(t)))dt,; J\_T(w)=(1)/(T)sum\_k: gamma\_k,gamma\_k+1in(T,2T](tau\_k+1-tau\_k) e\textasciicircum{}iwv\_k, v\_k=tau\_k+1/2,}
\expandafter\gdef\csname CGTip@thm:twist\endcsname{Theorem 8.12 (Transport identity) (p. 37): Uniformly for 0<w<=log T and T>= T\_0, I\_T(w)=J\_T(w)+(2i)/(T)sum\_kin I\_Tsin ((w)/(2L(tau\_k)))delta\_k e\textasciicircum{}iwtau\_k+O((1)/(T)). If moreover log 2<= w<=(1-varepsilon)log(T)/(2pi), then J\_T(w)<<\_varepsilon T\textasciicircum{}-1, so the first term may be dropped (the lower bound on w is needed: J\_T(w)->1 as w->0).}
\expandafter\gdef\csname CGTip@eq:twist_identity\endcsname{Eq. (8.17), p. 37: I\_T(w)=J\_T(w)+(2i)/(T)sum\_kin I\_Tsin ((w)/(2L(tau\_k)))delta\_k e\textasciicircum{}iwtau\_k+O((1)/(T)).}
\expandafter\gdef\csname CGTip@cor:weighted_response\endcsname{Corollary 8.13 (The weighted response at the clock points) (p. 37): Let p be a prime with theta=log p/L\_Tin[theta\_0,theta\_1]in(0,1), w=log p, and [ w\_k=(no.I\_T)/(TL(tau\_k))*(sin(w/2L(tau\_k)))/(sin(pi theta)),tilde R\textasciicircum{}tau\_p(I\_T)=-(2pisqrt(p))/(no.I\_T)sum\_kin I\_T w\_k s\_ksin(tau\_klog p). ] Then w\_k=1+O\_theta\_0,theta\_1(1/log T) and Re I\_T(log p)=(sin(pi theta))/(pisqrt(p)) tilde R\textasciicircum{}tau\_p(I\_T)...}
\expandafter\gdef\csname CGTip@eq:weighted_response\endcsname{Eq. (8.18), p. 37: Re I\_T(log p)=(sin(pi theta))/(pisqrt(p)) tilde R\textasciicircum{}tau\_p(I\_T)+O((1)/(T)).}
\expandafter\gdef\csname CGTip@cor:unweighted\endcsname{Corollary 8.14 (Removing the weights for slowly moving primes) (p. 38): Uniformly for primes p<= T\textasciicircum{}1/8, [ |tilde R\textasciicircum{}tau\_p(I\_T)-R\textasciicircum{}tau\_p(I\_T)|<<(1)/(log T)+(sqrt(p) log p log log T)/(log\textasciicircum{}2T)+(plog T)/(T), ] so that tilde R\textasciicircum{}tau\_p-R\textasciicircum{}tau\_p->0 uniformly for p<=(log T)\textasciicircum{}4-varepsilon; in particular, for p<=(log T)\textasciicircum{}2-varepsilon, Re I\_T(log p)=(sin(pi theta))/(pisqrt(p))(1+o(1)). For p=T\textasciicircum{}theta with...}
\expandafter\gdef\csname CGTip@prop:resonance\endcsname{Proposition 8.15 (Resonance of the clock sum) (p. 38): Let p be a prime with T/2pi <= p<= T/pi, and suppose min(2pi p-T,2T-2pi p)>=kappa T for a fixed kappa>0. Then [ J\_T(log p)=(2pisqrt(p))/(T)+O\_kappa (T\textasciicircum{}-3/5). ]}
\expandafter\gdef\csname CGTip@prop:pair_resonance\endcsname{Proposition 8.16 (The cross term of the primes at the clock points) (p. 39): Let I be a set of consecutive indices with tau\_kin[T\_1,T\_2] for kin I, T\_1<= T\_2<=2T\_1, T\_1>= T\_0, and let 2<= X<=(T\_1/2pi)\textasciicircum{}1/2. For integers n>=2 put d(n)=min(|2pi n-T\_1|,|2pi n-T\_2|). Then for every n<= T\_2/pi, sum\_kin Ie\textasciicircum{}itau\_klog n=-sqrt(n) log n 1\_[T\_1<2pi n<T\_2]+O(T\_1\textasciicircum{}2/5log T\_1+log T\_1min ((T\_1)/(d(n)),T\_1\textasciicircum{}1/...}
\expandafter\gdef\csname CGTip@eq:pair_resonance\endcsname{Eq. (8.19), p. 39: sum\_kin Ie\textasciicircum{}itau\_klog n=-sqrt(n) log n 1\_[T\_1<2pi n<T\_2]+O(T\_1\textasciicircum{}2/5log T\_1+log T\_1min ((T\_1)/(d(n)),T\_1\textasciicircum{}1/2)),}
\expandafter\gdef\csname CGTip@eq:cross_term_pairs\endcsname{Eq. (8.20), p. 39: Cr\textasciicircum{}(1)\_X(I)= (1)/(2pi\textasciicircum{}2no.I)sum\_substack p,q<= XT\_1<2pi pq<T\_2log(pq); +O((X\textasciicircum{}2log T\_1)/(no.Ilog\textasciicircum{}2X)+(sigma\_X\textasciicircum{}2)/(no.I)(T\_1\textasciicircum{}2/5log T\_1+log T\_1max\_p,q<= Xmin ((T\_1)/(d(pq)),T\_1\textasciicircum{}1/2))),}
\expandafter\gdef\csname CGTip@tab:twist\endcsname{Table 3 (p. 41): Transport identity (8.17) at T=2\textasciicircum{}19 (no.I\_T=977 804). I: the twisted integral, computed gap by gap; D: the displacement sum (2i)/(T)sum sin(*)delta\_kp\textasciicircum{}itau\_k; J: the clock sum; R\textasciicircum{}tau\_p and tilde R\textasciicircum{}tau\_p as in (6.1) and Corollary 8.13. The last two primes are in the resonance range of Proposition 8.15, where 2pisqrt(p)...}
\expandafter\gdef\csname CGTip@prop:almost_orthogonal\endcsname{Proposition 8.17 (Almost orthogonality at the clock points) (p. 41): Let 0<b<1, M=no.I\_T, v\_p(k)=sqrt(2/M) sin(tau\_klog p) for primes p<= T\textasciicircum{}b, and G=(< v\_p,v\_q>)\_p,q<= T\textasciicircum{}b with the inner product of R\textasciicircum{}M. Then ||G-I||\_rm op<<\_b(log T)\textasciicircum{}-1/4+T\textasciicircum{}b-1log T .}
\expandafter\gdef\csname CGTip@eq:almost_orthogonal\endcsname{Eq. (8.21), p. 41: ||G-I||\_rm op<<\_b(log T)\textasciicircum{}-1/4+T\textasciicircum{}b-1log T .}
\expandafter\gdef\csname CGTip@prop:spectral_energy\endcsname{Proposition 8.18 (Spectral energy and collective transfer) (p. 41): Let 0<b<1 and let (r\_k)\_kin I\_T be any real sequence. Then sum\_p<= T\textasciicircum{}b(1)/(p)(2pisqrt(p) < r\_ksin(tau\_klog p)>\_I\_T)\textasciicircum{}2<=2pi\textasciicircum{}2(1+o\_b(1))< r\_k\textasciicircum{}2>\_I\_T. In particular: item sum\_p<= T\textasciicircum{}b|R\textasciicircum{}tau\_p(I\_T)|\textasciicircum{}2/p<=2pi\textasciicircum{}2(1+o\_b(1))< s\_k\textasciicircum{}2>\_I\_T=log log T+O\_b(sqrt(log log T)). item For every sequence of points x\_k computed from the cl...}
\expandafter\gdef\csname CGTip@eq:spectral_energy\endcsname{Eq. (8.22), p. 42: sum\_p<= T\textasciicircum{}b(1)/(p)(2pisqrt(p) < r\_ksin(tau\_klog p)>\_I\_T)\textasciicircum{}2<=2pi\textasciicircum{}2(1+o\_b(1))< r\_k\textasciicircum{}2>\_I\_T.}
\expandafter\gdef\csname CGTip@eq:collective_transfer\endcsname{Eq. (8.23), p. 42: sum\_p<= T\textasciicircum{}b(|R\textasciicircum{}tau\_p(I\_T)- A\_p|\textasciicircum{}2)/(p)<=2pi\textasciicircum{}2(1+o\_b(1))< r\_k\textasciicircum{}2>\_I\_T, sum\_substack p<= T\textasciicircum{}b||R\textasciicircum{}tau\_p- A\_p|>varepsilon(1)/(p)<=(2pi\textasciicircum{}2+o\_b(1))/(varepsilon\textasciicircum{}2)< r\_k\textasciicircum{}2>\_I\_T.}
\expandafter\gdef\csname CGTip@lem:sorted\endcsname{Lemma 8.19 (Sorted Newton steps) (p. 42): Let gamma\_1<=...<=gamma\_M and let z\_1,...,z\_M be arbitrary reals, with z\textasciicircum{}increasing to the same values in increasing order. Then sum\_k|z\textasciicircum{}increasing to\_k-gamma\_k|\textasciicircum{}q<=sum\_k|z\_k-gamma\_k|\textasciicircum{}q for every q>=1, and z\textasciicircum{}increasing to\_k=z\_k as multisets. Consequently, replacing x\textasciicircum{}(j)\_k by its sorted version at any level of (8.1)...}
\expandafter\gdef\csname CGTip@sec:envelope\endcsname{Section 9: Theorem E: the pointwise envelope (p. 43)}
\expandafter\gdef\csname CGTip@prop:envelope_equivalence\endcsname{Proposition 9.1 (The envelope of the clock is the envelope of S) (p. 43): (a) limsup\_k->inf|tau\_k-gamma\_k|=2pi a. (b) If X\_k>=2 and sigma\_X\_k=o(log tau\_k), in particular if X\_k<=(log tau\_k)\textasciicircum{}2-varepsilon, then limsup\_k->inf|tau\_k\textasciicircum{}X\_k-gamma\_k|=2pi a.}
\expandafter\gdef\csname CGTip@prop:finite_envelope\endcsname{Proposition 9.2 (Finite-height envelope) (p. 43): For every k>=1 let E\_k=BWS(gamma\_k)-(1)/(2)+|varpi(gamma\_k)|/pi, and E\_k=2.5167-(1)/(2)+|varpi(gamma\_k)|/pi if gamma\_k<=30 610 046 000. Then |s\_k|<= E\_k and -(E\_k)/(L(tau\_k))<=delta\_k<=(E\_k)/(L(gamma\_k)) .}
\expandafter\gdef\csname CGTip@eq:finite_envelope\endcsname{Eq. (9.1), p. 43: -(E\_k)/(L(tau\_k))<=delta\_k<=(E\_k)/(L(gamma\_k)) .}
\expandafter\gdef\csname CGTip@thm:envelope\endcsname{Theorem 9.3 (Theorem E) (p. 44): Let e be continuous on [(1)/(2),1] with e(1)=0, such that for every varepsilon>0, zeta(sigma+it)<<\_varepsilon|t|\textasciicircum{} e(sigma)+varepsilon uniformly for (1)/(2)<=sigma <=1 and |t|>=2, and put e(sigma)=(1)/(2)-sigma+ e(1-sigma) for 0<=sigma<(1)/(2). Then a<=(1)/(2pi log 2)int\_0\textasciicircum{}1(e(sigma))/(sqrt(sigma(2-sigma))) dsigma ....}
\expandafter\gdef\csname CGTip@eq:envelope_formula\endcsname{Eq. (9.2), p. 44: a<=(1)/(2pi log 2)int\_0\textasciicircum{}1(e(sigma))/(sqrt(sigma(2-sigma))) dsigma .}
\expandafter\gdef\csname CGTip@eq:envelope_value\endcsname{Eq. (9.3), p. 44: limsup\_k->inf|tau\_k-gamma\_k|=2pi a<=0.6083129...<0.60832 .}
\expandafter\gdef\csname CGTip@rem:envelope_inputs\endcsname{Remark 9.4 (p. 45): (a) With the explicit exponents used by Bellotti and Wong ((27)/(164) at sigma=(1)/(2), Yang's lines sigma\_j=1-j/(2\textasciicircum{}j-2) with exponents 1/(2\textasciicircum{}j-2), 4<= j<=9 [ref]) and linear interpolation between sigma=0 and sigma=(1)/(2), as in [ref, S 3.1], the right side of (9.2) is 0.10074, the limit of their C\_1=0.10076 as thei...}
\expandafter\gdef\csname CGTip@prop:explicit_S\endcsname{Proposition 9.5 (An explicit bound valid for all T>= e) (p. 45): For T>= e, |S(T)|<=0.09783log T+0.722log log T+13,; |N(T)-(T)/(2pi)log(T)/(2pi e)|<=0.09783log T+0.722log log T+14 .}
\expandafter\gdef\csname CGTip@eq:explicit_S\endcsname{Eq. (9.4), p. 45: |S(T)|<=0.09783log T+0.722log log T+13,; |N(T)-(T)/(2pi)log(T)/(2pi e)|<=0.09783log T+0.722log log T+14 .}
\expandafter\gdef\csname CGTip@sec:computations\endcsname{Section 10: Computations (p. 46)}
\expandafter\gdef\csname CGTip@subsec:data\endcsname{Section 10.1: Data and methods (p. 46)}
\expandafter\gdef\csname CGTip@subsec:clock_data\endcsname{Section 10.2: The clock and the ordinates (p. 46)}
\expandafter\gdef\csname CGTip@fig:displacement\endcsname{Figure 1 (p. 47): The displacement delta\_k=tau\_k-gamma\_k of the RvM clock and the residual tau\_k\textasciicircum{}PP-gamma\_k of the one-step clock with PP=2,3,5,7, for k<=2000; the correction S\_PP(tau\_k)/L(tau\_k) is capped at +- B with B=0.6 (saturation, Section 10.5).}
\expandafter\gdef\csname CGTip@subsec:direct_defects\endcsname{Section 10.3: Direct defects (Theorem B) (p. 46)}
\expandafter\gdef\csname CGTip@tab:direct\endcsname{Table 4 (p. 47): Mean squares of the defect u\_k(X) and of the one-step residual r\_k(X) over T<gamma\_k<=2T, X=T\textasciicircum{}c.}
\expandafter\gdef\csname CGTip@fig:direct_defects\endcsname{Figure 2 (p. 48): Mean square of u\_k(X) (filled markers) and of r\_k(X) (open markers) over T<gamma\_k<=2T against log(1/c), X=T\textasciicircum{}c, for the four heights of Table 4; dashed: (1)/(2pi\textasciicircum{}2)log(1/c).}
\expandafter\gdef\csname CGTip@subsec:accounting_data\endcsname{Section 10.4: Exact accounting (Theorem C) (p. 46)}
\expandafter\gdef\csname CGTip@tab:accounting\endcsname{Table 5 (p. 48): Exact decomposition (6.5) of the one-step gain into the diagonal term D\_X and the cross term Cr\_X.}
\expandafter\gdef\csname CGTip@fig:accounting\endcsname{Figure 3 (p. 49): Gain of the one-step clock (filled markers) and cross term Cr\_X (open markers) against the cutoff X, for five windows of ordinals; dashed: v\_X.}
\expandafter\gdef\csname CGTip@subsec:one_step_cutoff\endcsname{Section 10.5: The one-step clock against the cutoff (p. 50)}
\expandafter\gdef\csname CGTip@tab:methods\endcsname{Table 6 (p. 50): Root mean square, maximum of |tau\_k\textasciicircum{}X-gamma\_k| and number of residuals >=(1)/(2) over all K ordinates, and the same quantities in the top window (10\textasciicircum{}6,K].}
\expandafter\gdef\csname CGTip@fig:rms_cutoff\endcsname{Figure 4 (p. 51): Root mean square of tau\_k\textasciicircum{}X-gamma\_k (no saturation) against the cutoff X, by window of ordinals; dots mark the minimum in each window.}
\expandafter\gdef\csname CGTip@fig:max_window\endcsname{Figure 5 (p. 51): Maximum of |tau\_k\textasciicircum{}X-gamma\_k| in each window of ordinals for the five rules of Table 6, logarithmic scale; dotted: (1)/(2). The value 3.247 of the fixed set p<=1999 in the first window occurs at k=1.}
\expandafter\gdef\csname CGTip@subsec:transport_data\endcsname{Section 10.6: Transport and the Newton term (Theorem D) (p. 50)}
\expandafter\gdef\csname CGTip@tab:transport\endcsname{Table 7 (p. 51): Transport term and Newton term of the one-step clock over 2\textasciicircum{}19<gamma\_k<=2\textasciicircum{}20, X=T\textasciicircum{}c: mean squares, the correlation of u\_k(X) with t\_k(X), the root mean square of S\_X'(tau\_k)/L(tau\_k), and the proportion of indices with F\_X'(tau\_k)<0.}
\expandafter\gdef\csname CGTip@subsec:onset\endcsname{Section 10.7: Onset of the cross term (p. 52)}
\expandafter\gdef\csname CGTip@tab:onset\endcsname{Table 8 (p. 52): Cross term at the clock points on (10\textasciicircum{}6,K]; (tau\_min/2pi)\textasciicircum{}1/2=309.1.}
\expandafter\gdef\csname CGTip@sec:responses\endcsname{Section 11: Prime responses at the zeros and at the clock points (p. 52)}
\expandafter\gdef\csname CGTip@subsec:landau_line\endcsname{Section 11.1: Fixed and slowly moving frequencies (p. 52)}
\expandafter\gdef\csname CGTip@thm:landau_line\endcsname{Theorem 11.1 (Landau's formula on the critical line) (p. 52): Uniformly for t>=2 and integers 2<= n<= t\textasciicircum{}1/8, sum\_0<gamma <= tn\textasciicircum{}igamma=-(t)/(2pi)(Lambda(n))/(sqrt(n))+O((tlog\textasciicircum{}2n)/(log t)+sqrt(n)log\textasciicircum{}2t), zeros counted with multiplicity. For fixed n the error is O(t/log t).}
\expandafter\gdef\csname CGTip@eq:landau_line\endcsname{Eq. (11.1), p. 52: sum\_0<gamma <= tn\textasciicircum{}igamma=-(t)/(2pi)(Lambda(n))/(sqrt(n))+O((tlog\textasciicircum{}2n)/(log t)+sqrt(n)log\textasciicircum{}2t),}
\expandafter\gdef\csname CGTip@lem:clock_equidistribution\endcsname{Lemma 11.2 (p. 53): For fixed real omega !=0, sum\_kin I\_Te\textasciicircum{}iomega tau\_k<<\_omega log T. Hence for fixed integers n\_1,n\_2>=2, <sin(tau\_klog n\_1)sin(tau\_klog n\_2)>\_I\_T->(1)/(2) 1\_n\_1=n\_2.}
\expandafter\gdef\csname CGTip@thm:projection\endcsname{Theorem 11.3 (p. 53): Uniformly for T>= T\_0 and integers 2<= n<= T\textasciicircum{}1/8, < s\_ksin(gamma\_klog n)>\_I\_T=(a\_n)/(2)+O((log n log log T)/(log T)+(sqrt(n)log\textasciicircum{}2T)/(Tlog n)), and the same with tau\_k in place of gamma\_k. Hence for primes p<= T\textasciicircum{}1/8, R\textasciicircum{}gamma\_p(I\_T) and R\textasciicircum{}tau\_p(I\_T) are 1+O((sqrt(p)log plog log T)/(log T)+(plog\textasciicircum{}2T)/(Tlog p)); they ten...}
\expandafter\gdef\csname CGTip@eq:projection_uniform\endcsname{Eq. (11.2), p. 53: < s\_ksin(gamma\_klog n)>\_I\_T=(a\_n)/(2)+O((log n log log T)/(log T)+(sqrt(n)log\textasciicircum{}2T)/(Tlog n)),}
\expandafter\gdef\csname CGTip@eq:projection_Phi\endcsname{Eq. (11.3), p. 53: H\_n(t)=[(S(y)n\textasciicircum{}iy)/(il\_n)]\_14\textasciicircum{}t-(1)/(il\_n)sum\_14<gamma\_k<= tn\textasciicircum{}igamma\_k+(1)/(il\_n)int\_14\textasciicircum{}tn\textasciicircum{}iy(L+(varpi')/(pi))dy=(i)/(l\_n)sum\_0<gamma\_k<= tn\textasciicircum{}igamma\_k+O(log t),}
\expandafter\gdef\csname CGTip@eq:projection_partial\endcsname{Eq. (11.4), p. 54: int\_T\textasciicircum{}2TS(t)n\textasciicircum{}itL(t) dt=[ H\_nL ]\_T\textasciicircum{}2T-int\_T\textasciicircum{}2T(H\_n(t))/(2pi t) dt=-i(Lambda(n))/(2pisqrt(n) l\_n)int\_T\textasciicircum{}2TL dt+O(Tl\_n+(sqrt(n)log\textasciicircum{}3T)/(l\_n)).}
\expandafter\gdef\csname CGTip@prop:cross_term\endcsname{Proposition 11.4 (No cross term below T\textasciicircum{}1/4) (p. 54): For varepsilon>0 there is T\_0(varepsilon) such that for T>= T\_0(varepsilon) and 2<= X<= T\textasciicircum{}1/4-varepsilon, |Cr\_X(I\_T)|<< T\textasciicircum{}-varepsilon/2sqrt(log log T).}
\expandafter\gdef\csname CGTip@cor:gain\endcsname{Corollary 11.5 (Gain at the clock points) (p. 54): (a) For a finite set A of prime powers and lambda in R, < s\_k\textasciicircum{}2-(s\_k-lambda S\_A(tau\_k))\textasciicircum{}2>\_I\_T->(2lambda-lambda\textasciicircum{}2)v\_A. (b) For fixed varepsilon>0 and lambda, uniformly for 2<= X<=(log T)\textasciicircum{}2-varepsilon, with full local factors, [ < s\_k\textasciicircum{}2-(s\_k-lambda S\_X(tau\_k))\textasciicircum{}2>\_I\_T=(2lambda-lambda\textasciicircum{}2)v\_X+O\_lambda ((sqrt(X)log log T)...}
\expandafter\gdef\csname CGTip@fig:gain_fixed\endcsname{Figure 6 (p. 55): Mean-square gain < s\_k\textasciicircum{}2>-<(s\_k-S\_PP(tau\_k))\textasciicircum{}2> of two fixed prime sets PP (full local factors) in each window of ordinals, and the constants v\_PP (dashed) of Corollary 11.5.}
\expandafter\gdef\csname CGTip@subsec:gamma_response\endcsname{Section 11.2: The response at the ordinates and its Fourier reduction (p. 55)}
\expandafter\gdef\csname CGTip@eq:ATp_STp\endcsname{Eq. (11.5), p. 55: A\_T,p=-(2pisqrt(p))/(no.I\_T)int\_T\textasciicircum{}2TG(t)sin(tlog p)L(t) dt, S\_T,p=(1)/(T)int\_T\textasciicircum{}2TS(t)\textasciicircum{}2p\textasciicircum{}it dt.}
\expandafter\gdef\csname CGTip@prop:gamma_response_exact\endcsname{Proposition 11.6 (Unconditional separation of one- and two-point terms) (p. 55): R\textasciicircum{}gamma\_p(I\_T)= A\_T,p+(pisqrt(p) Tlog p)/(no.I\_T)Re S\_T,p-(pisqrt(p))/(no.I\_T)[G\textasciicircum{}2sin(tlog p)]\_T\textasciicircum{}2T; +(pisqrt(p)log p)/(no.I\_T)int\_T\textasciicircum{}2T(G\textasciicircum{}2-S\textasciicircum{}2)cos(tlog p) dt. If theta=log p/L\_T stays in a compact subset of (0,1), the last two terms are o(1) and R\textasciicircum{}gamma\_p(I\_T)= A\_T,p+2pi\textasciicircum{}2thetasqrt(p) (1+O(1/L\_T))Re S\_T,p+o(1).}
\expandafter\gdef\csname CGTip@eq:gamma_response_exact\endcsname{Eq. (11.6), p. 55: R\textasciicircum{}gamma\_p(I\_T)= A\_T,p+(pisqrt(p) Tlog p)/(no.I\_T)Re S\_T,p-(pisqrt(p))/(no.I\_T)[G\textasciicircum{}2sin(tlog p)]\_T\textasciicircum{}2T; +(pisqrt(p)log p)/(no.I\_T)int\_T\textasciicircum{}2T(G\textasciicircum{}2-S\textasciicircum{}2)cos(tlog p) dt.}
\expandafter\gdef\csname CGTip@eq:gamma_response_scaled\endcsname{Eq. (11.7), p. 55: R\textasciicircum{}gamma\_p(I\_T)= A\_T,p+2pi\textasciicircum{}2thetasqrt(p) (1+O(1/L\_T))Re S\_T,p+o(1).}
\expandafter\gdef\csname CGTip@eq:measure_def\endcsname{Eq. (11.8), p. 56: Z\_T(omega)=int\_R e\textasciicircum{}iomega t d m\_T(t), hat f\_T(omega)=int\_R f\_T(t)e\textasciicircum{}iomega tdt .}
\expandafter\gdef\csname CGTip@thm:Fourier_exact\endcsname{Theorem 11.7 (Unconditional shifted two-point representation) (p. 56): For non-ordinates T,2T, Z\_T(omega)=sum\_T<gamma<2Tnu(gamma)e\textasciicircum{}iomega gamma-int\_T\textasciicircum{}2T(L(t)+(varpi'(t))/(pi))e\textasciicircum{}iomega t dt+S(T)e\textasciicircum{}iomega T-S(2T)e\textasciicircum{}2iomega T, Z\_T(omega)=-iomega hat f\_T(omega), so that Z\_T(0)=0 and Z\_T(omega)/omega extends continuously by -ihat f\_T; and for every real omega\_0 int\_T\textasciicircum{}2TS(t)\textasciicircum{}2e\textasciicircum{}iomega\_0t dt=-(...}
\expandafter\gdef\csname CGTip@eq:measure_expansion\endcsname{Eq. (11.9), p. 56: Z\_T(omega)=sum\_T<gamma<2Tnu(gamma)e\textasciicircum{}iomega gamma-int\_T\textasciicircum{}2T(L(t)+(varpi'(t))/(pi))e\textasciicircum{}iomega t dt+S(T)e\textasciicircum{}iomega T-S(2T)e\textasciicircum{}2iomega T,}
\expandafter\gdef\csname CGTip@eq:Fourier_convolution\endcsname{Eq. (11.10), p. 56: int\_T\textasciicircum{}2TS(t)\textasciicircum{}2e\textasciicircum{}iomega\_0t dt=-(1)/(2pi)int\_R(Z\_T(omega)Z\_T(omega\_0-omega))/(omega(omega\_0-omega)) domega,}
\expandafter\gdef\csname CGTip@eq:BT_scaled\endcsname{Eq. (11.11), p. 56: B\_T(theta)=int\_R(Z\_T(alpha L\_T)Z\_T((theta-alpha)L\_T))/(alpha(theta-alpha)) dalpha,so that S\_T,p=-(B\_T(theta))/(2pi TL\_T).}
\expandafter\gdef\csname CGTip@subsec:kernel\endcsname{Section 11.3: The kernel K\_theta (p. 56)}
\expandafter\gdef\csname CGTip@eq:kernel_piecewise\endcsname{Eq. (11.12), p. 56: K\_theta(alpha)= theta, alpha <=-1,; theta-1-alpha, -1<=alpha <=theta-1,; 0, theta-1<=alpha <=1,; alpha-1, 1<=alpha <=1+theta,; theta, alpha >=1+theta,}
\expandafter\gdef\csname CGTip@eq:kernel_intersections\endcsname{Eq. (11.13), p. 57: K\_theta(alpha)=1-h(alpha)-h(theta-alpha)-h(theta)+2max (1-max(|alpha|,|theta-alpha|,theta),0),}
\expandafter\gdef\csname CGTip@prop:kernel_integral\endcsname{Proposition 11.8 (Integral and support of the kernel) (p. 57): For 0<theta<1, int\_R(K\_theta(alpha))/(alpha(theta-alpha)) dalpha=-(2Phi(theta))/(theta), the integral being absolutely convergent, and K\_theta(alpha)=0 whenever |alpha|+|theta-alpha|+theta<2.}
\expandafter\gdef\csname CGTip@eq:kernel_integral\endcsname{Eq. (11.14), p. 57: int\_R(K\_theta(alpha))/(alpha(theta-alpha)) dalpha=-(2Phi(theta))/(theta),}
\expandafter\gdef\csname CGTip@prop:PY_square\endcsname{Proposition 11.9 (The square of the prime polynomial has no prime frequency) (p. 57): For Y>=2, T>=2 and every prime p, |(1)/(T)int\_T\textasciicircum{}2TS\textasciicircum{}rm pr\_Y(t)\textasciicircum{}2p\textasciicircum{}it dt|<< Y\textasciicircum{}3/T, with an absolute constant. In particular, with E\_0=S-S\textasciicircum{}rm pr\_Y\_0, uniformly for primes p<= T\textasciicircum{}theta\_1 with fixed theta\_1<1, [ S\_T,p=(1)/(T)int\_T\textasciicircum{}2T(2S\textasciicircum{}rm pr\_Y\_0(t)E\_0(t)+E\_0(t)\textasciicircum{}2)p\textasciicircum{}it dt+o(p\textasciicircum{}-1/2). ]}
\expandafter\gdef\csname CGTip@subsec:newton_spectrum\endcsname{Section 11.4: The arithmetic of the Newton nonlinearity (p. 57)}
\expandafter\gdef\csname CGTip@prop:spectral_barrier\endcsname{Proposition 11.10 (Spectral barrier at the clock points) (p. 58): Let 0<theta<1, let p=p(T) be a prime with log p/log T->theta, let epsilon\_T->0, and let [ Q\_T(t)=sum\_a,b<= T\textasciicircum{}epsilon\_Tc\_a,b e\textasciicircum{}itlog(a/b),sum\_a,b|c\_a,b|<= T\textasciicircum{}o(1), ] with integers a,b>=1 and complex coefficients. Then [ (sqrt(p))/(M)sum\_kin I\_TQ\_T(tau\_k) p\textasciicircum{} itau\_k=o(1). ]}
\expandafter\gdef\csname CGTip@lem:bessel\endcsname{Lemma 11.11 (Bessel expansion of a Newton phase) (p. 58): Let X>=2 and lambda in R, and let P\_X=S\textasciicircum{}(1)\_X be the prime part of S\_X, [ P\_X(t)=-(1)/(pi)sum\_q<= X(sin(tlog q))/(sqrt(q)). ] With z\_q=lambda/(pisqrt(q)) and integer vectors j=(j\_q)\_q<= X indexed by the primes q<= X, e\textasciicircum{}-ilambda P\_X(t)=sum\_ jA\_ j e\textasciicircum{}itLambda\_ j, A\_ j=prod\_q<= XJ\_j\_q(z\_q),Lambda\_ j=sum\_q<= Xj\_qlog q, t...}
\expandafter\gdef\csname CGTip@eq:bessel_expansion\endcsname{Eq. (11.15), p. 58: e\textasciicircum{}-ilambda P\_X(t)=sum\_ jA\_ j e\textasciicircum{}itLambda\_ j, A\_ j=prod\_q<= XJ\_j\_q(z\_q),Lambda\_ j=sum\_q<= Xj\_qlog q,}
\expandafter\gdef\csname CGTip@eq:bessel_moments\endcsname{Eq. (11.16), p. 58: sum\_ jA\_ j\textasciicircum{}2=1,sum\_ jA\_ j\textasciicircum{}2Lambda\_ j=0,sum\_ jA\_ j\textasciicircum{}2Lambda\_ j\textasciicircum{}2=(lambda\textasciicircum{}2)/(2pi\textasciicircum{}2)sum\_q<= X(log\textasciicircum{}2q)/(q),}
\expandafter\gdef\csname CGTip@eq:bessel_tail\endcsname{Eq. (11.17), p. 58: sum\_D(j)>= DA\_ j\textasciicircum{}2<= e\textasciicircum{}-sDexp ((C\_slambda\textasciicircum{}2)/(pi\textasciicircum{}2)sum\_q<= X(1)/(q)), C\_s=e\textasciicircum{}sexp(e\textasciicircum{}s/4+1/2).}
\expandafter\gdef\csname CGTip@prop:newton_spectrum\endcsname{Proposition 11.12 (Spectral concentration and dispersion of the Newton phases) (p. 59): Let X\_1=T\textasciicircum{}c\_1, c\_1=H\_0\textasciicircum{}-1/2, and for a prime power n and t>2pi let lambda\_n(t)=log n/L(t), so that sin(x\textasciicircum{}(1)\_klog n)=Im (e\textasciicircum{}itau\_klog ne\textasciicircum{}-ilambda\_n(tau\_k)S\_X\_1(tau\_k)); the prime-power part S\_X\_1-P\_X\_1 contributes to e\textasciicircum{}-ilambda S\_X\_1 a further factor of the form (11.15), with z\_p\textasciicircum{}m=lambda/(pi mp\textasciicircum{}m/2) at the frequenci...}
\expandafter\gdef\csname CGTip@eq:newton_dispersion\endcsname{Eq. (11.18), p. 59: sum\_|Lambda\_ j|>=theta log TA\_ j\textasciicircum{}2<=(c\_1\textasciicircum{}2c\_2\textasciicircum{}2)/(theta\textasciicircum{}2)(1+o(1))=(1+o(1))/(theta\textasciicircum{}2H\_0H\_1).}
\expandafter\gdef\csname CGTip@prop:orthogonality\endcsname{Proposition 11.13 (Exact orthogonality for a fixed cutoff) (p. 60): Let X>=2 be fixed, let F be a continuous function on the torus T\textasciicircum{}no.q<= X, and let F\_X(t)=F((tlog q)\_q<= X). Then for every prime p>X, [ lim\_U->inf(1)/(U)int\_U\textasciicircum{}2Ue\textasciicircum{}itlog pF\_X(t) dt=0 . ]}
\expandafter\gdef\csname CGTip@eq:transfer_estimate\endcsname{Eq. (11.19), p. 60: sup\_T\textasciicircum{}a<= p<= T\textasciicircum{}bsqrt(p) |(1)/(M)sum\_kin I\_TL(tau\_k)(x\textasciicircum{}(m)\_k-gamma\_k)e\textasciicircum{}itau\_klog p|->0(0<a<b<1),}
\expandafter\gdef\csname CGTip@eq:computable_parts\endcsname{Eq. (11.20), p. 60: A\textasciicircum{}(m)\_p =-2pisqrt(p)<L(tau\_k)(tau\_k-x\textasciicircum{}(m)\_k)sin(tau\_klog p)>\_I\_T,; E\textasciicircum{}(m)\_p =-2pisqrt(p)<L(tau\_k)(x\textasciicircum{}(m)\_k-gamma\_k)sin(tau\_klog p)>\_I\_T,; D\textasciicircum{}(m)\_p =-2pisqrt(p)< S\_X\_m(tau\_k)sin(tau\_klog p)>\_I\_T,; B\textasciicircum{}(m)\_p =-2pisqrt(p)< B\textasciicircum{}(m)\_ksin(tau\_klog p)>\_I\_T, B\textasciicircum{}(m)\_k=S\_X\_m(x\textasciicircum{}(m-1)\_k)-S\_X\_m(tau\_k),}
\expandafter\gdef\csname CGTip@thm:computable_response\endcsname{Theorem 11.14 (Computable and residual responses) (p. 60): Let X\_m=T\textasciicircum{}c\_m with c\_m<=(1)/(2) and x\textasciicircum{}(m)\_k as in (8.1). Then: item R\textasciicircum{}tau\_p(I\_T)= A\textasciicircum{}(m)\_p+ E\textasciicircum{}(m)\_p+O(sqrt(p)/T) and A\textasciicircum{}(m)\_p= D\textasciicircum{}(m)\_p+ B\textasciicircum{}(m)\_p+O(sqrt(p) T\textasciicircum{}c\_m-1), for every m>=1 and every prime p; B\textasciicircum{}(1)\_p=0. item For varepsilon>0, uniformly for primes p with pX\_m<= T\textasciicircum{}1-varepsilon, D\textasciicircum{}(m)\_p= 1\_[p<= X\_m]+O\_varepsilon ((...}
\expandafter\gdef\csname CGTip@eq:diagonal_response\endcsname{Eq. (11.21), p. 61: D\textasciicircum{}(m)\_p= 1\_[p<= X\_m]+O\_varepsilon (((p+X\_m)log T)/(T)).}
\expandafter\gdef\csname CGTip@eq:moving_response_split\endcsname{Eq. (11.22), p. 61: R\textasciicircum{}tau\_p(I\_T)= B\textasciicircum{}(m)\_p+ E\textasciicircum{}(m)\_p+o(1):}
\expandafter\gdef\csname CGTip@tab:computable_response\endcsname{Table 9 (p. 62): Computable and residual responses at T=2\textasciicircum{}19 (no.I\_T=977 804) for the three cutoff sets of Table 2: R\textasciicircum{}tau\_p and, for m=1,2,3, the diagonal part D\textasciicircum{}(m)\_p, the transport part B\textasciicircum{}(m)\_p and the residual E\textasciicircum{}(m)\_p of (11.20). A dagger marks pX\_m>=83 443 (pair resonance); prime powers n=q\textasciicircum{}j with q<= X\_m and the width of the stat...}
\expandafter\gdef\csname CGTip@subsec:conditional_law\endcsname{Section 11.5: The moving-prime law under Hypotheses U and W (p. 62)}
\expandafter\gdef\csname CGTip@hyp:uniform_prime\endcsname{Hypothesis 11.15 (Hypothesis U: uniform one-particle input) (p. 62): For every 0<theta\_0<theta\_1<1, uniformly for primes with theta\_0<=log p/L\_T<=theta\_1, A\_T,p=1+o(1).}
\expandafter\gdef\csname CGTip@hyp:weighted_bridge\endcsname{Hypothesis 11.16 (Hypothesis W: weighted arithmetic bridge) (p. 62): Uniformly for primes with theta in any compact subset of (0,1), int\_R(Z\_T(alpha L\_T)Z\_T((theta-alpha)L\_T)+ N\_Tp\textasciicircum{}-1/2 K\_theta(alpha))/(alpha(theta-alpha)) dalpha=o((TL\_T)/(sqrt(p))), with the continuous extensions of Z\_T(omega)/omega at omega=0.}
\expandafter\gdef\csname CGTip@eq:weighted_arithmetic_bridge\endcsname{Eq. (11.23), p. 62: int\_R(Z\_T(alpha L\_T)Z\_T((theta-alpha)L\_T)+ N\_Tp\textasciicircum{}-1/2 K\_theta(alpha))/(alpha(theta-alpha)) dalpha=o((TL\_T)/(sqrt(p))),}
\expandafter\gdef\csname CGTip@prop:conditional_universal\endcsname{Proposition 11.17 (Moving-prime law under Hypotheses U and W) (p. 63): Under Hypothesis W, S\_T,p=-(1)/(2pi\textasciicircum{}2sqrt(p))(Phi(theta))/(theta)+o(p\textasciicircum{}-1/2). Under Hypotheses U and W, or U and the real part of (11.24), uniformly for theta in compact subsets of (0,1), R\textasciicircum{}gamma\_p(I\_T)=1-Phi(theta)+o(1). Conversely, under Hypothesis U, (11.25) is equivalent to the real part of (11.24). Nothing is as...}
\expandafter\gdef\csname CGTip@eq:CTp_target\endcsname{Eq. (11.24), p. 63: S\_T,p=-(1)/(2pi\textasciicircum{}2sqrt(p))(Phi(theta))/(theta)+o(p\textasciicircum{}-1/2).}
\expandafter\gdef\csname CGTip@eq:conditional_universal\endcsname{Eq. (11.25), p. 63: R\textasciicircum{}gamma\_p(I\_T)=1-Phi(theta)+o(1).}
\expandafter\gdef\csname CGTip@rem:scale_gap\endcsname{Remark 11.18 (p. 63): Phi(theta)=2sum\_j>=1(theta\textasciicircum{}2j)/((2j-1)2j)=theta\textasciicircum{}2+(theta\textasciicircum{}4)/(6)+(theta\textasciicircum{}6)/(15)+..., so 1-theta\textasciicircum{}2 is only the quadratic truncation of (11.25), and 1-Phi(theta) vanishes at theta=0.9036.... Hypotheses U and W are not consequences of Theorem 11.1, which reaches only theta ->0. By Proposition 11.9, (11.24) is a statemen...}
\expandafter\gdef\csname CGTip@sec:heuristic\endcsname{Section 12: Heuristic support and numerical evidence (p. 63)}
\expandafter\gdef\csname CGTip@subsec:det_model\endcsname{Section 12.1: A determinantal model whose connected kernel is K\_theta (p. 63)}
\expandafter\gdef\csname CGTip@eq:real_tilt\endcsname{Eq. (12.1), p. 63: d P\_iota\_1,iota\_2=(exp -iota\_1Re Xi\_-theta-iota\_2Im Xi\_-theta)/(Eexp -iota\_1Re Xi\_-theta-iota\_2Im Xi\_-theta) d P\_0,0}
\expandafter\gdef\csname CGTip@thm:sine_kernel\endcsname{Theorem 12.1 (Connected kernel of the model) (p. 63): Let lattice frequencies tend to alpha,theta-alpha,-theta with 0<theta<1. Then R\_0\textasciicircum{}-1cum\_3(Xi\_alpha,Xi\_theta-alpha,Xi\_-theta)-> K\_theta(alpha) and R\_0\textasciicircum{}-1cum\_2(Xi\_theta,Xi\_-theta)->theta.}
\expandafter\gdef\csname CGTip@eq:third_cumulant_trace\endcsname{Eq. (12.2), p. 64: cum\_3= Tr(m\_1 m\_2 m\_3Pi\_ n)-Tr(m\_1 m\_2Pi\_ n m\_3Pi\_ n)-Tr(m\_1 m\_3Pi\_ n m\_2Pi\_ n)-Tr(m\_2 m\_3Pi\_ n m\_1Pi\_ n); +Tr(m\_1Pi\_ n m\_2Pi\_ n m\_3Pi\_ n)+Tr(m\_1Pi\_ n m\_3Pi\_ n m\_2Pi\_ n),}
\expandafter\gdef\csname CGTip@subsec:numerics_law\endcsname{Section 12.2: Numerical evidence (p. 64)}
\expandafter\gdef\csname CGTip@conj:law\endcsname{Conjecture 12.2 (p. 64): For primes with theta=log p/L\_T in compact subsets of (0,1), R\textasciicircum{}gamma\_p(I\_T)=1-Phi(theta)+o(1).}
\expandafter\gdef\csname CGTip@tab:response_errors\endcsname{Table 10 (p. 64): Errors of R\textasciicircum{}gamma\_p-(1-Phi(theta)) over the primes p<=20 000 with theta<1.}
\expandafter\gdef\csname CGTip@tab:response_bins\endcsname{Table 11 (p. 65): Bin means of R\textasciicircum{}gamma\_p and R\textasciicircum{}tau\_p in the two top windows; ``law'' is the mean of 1-Phi(theta) over the same primes, s.e. the standard error of the mean difference.}
\expandafter\gdef\csname CGTip@fig:response_law\endcsname{Figure 7 (p. 66): Bin means of R\textasciicircum{}gamma\_p (bars: descriptive standard error, see the text) against 1-Phi(theta) (solid) and 1-theta\textasciicircum{}2 (dashed), for five windows of ordinals.}
\expandafter\gdef\csname CGTip@fig:response_clock\endcsname{Figure 8 (p. 66): Bin means of the response at the clock points, R\textasciicircum{}tau\_p (squares), and at the zeros, R\textasciicircum{}gamma\_p (circles), for three windows of ordinals; solid: 1-Phi(theta).}
\expandafter\gdef\csname CGTip@sec:conclusions\endcsname{Section 13: Conclusions (p. 67)}
\expandafter\gdef\csname CGTip@subsec:pending\endcsname{Section 13.3: Status of the proofs and pending verification (p. 69)}
\expandafter\gdef\csname CGTip@subsec:missing\endcsname{Section 13.4: What is missing for a complete prediction of the zeros from the clock (p. 70)}
\expandafter\gdef\csname CGTip@q:constant\endcsname{Question 13.1 (p. 71): For fixed cin(0,(1)/(2)], do the bounded means of u\_k(T\textasciicircum{}c)\textasciicircum{}2 converge, and to (2pi\textasciicircum{}2)\textasciicircum{}-1log(1/c) plus which constant?}
\expandafter\gdef\csname CGTip@q:fixed_c\endcsname{Question 13.2 (p. 71): Does < r\_k(T\textasciicircum{}c)\textasciicircum{}2>\_I\_T stay bounded for fixed c, that is, does the transport cancel against the defect as in Table 7? Is the error of the corrected clock after log\textasciicircum{}*T+O(1) steps bounded in units of zeros for almost every index, and not only in mean square? Can Theorems 5.1 and 7.3 be extended to c>(1)/(2)?}
\expandafter\gdef\csname CGTip@q:prime_bridge\endcsname{Question 13.3 (p. 71): Can Hypotheses U and W be proved for p=T\textasciicircum{}theta, that is, can the terms of Proposition 11.9 that contain Selberg's remainder be evaluated to o(p\textasciicircum{}-1/2)? What is the finite-height correction to 1-Phi(theta) seen in Table 11?}
\expandafter\gdef\csname CGTip@app:programs\endcsname{Section Appendix ~A: Programs (p. 74)}
\makeatother
